\documentclass[11pt,a4paper]{article}
\usepackage{fontspec}
\usepackage{amsmath,amssymb,amsthm,mathtools}
\usepackage[margin=24mm]{geometry}
\usepackage{xspace}
\usepackage{xcolor,booktabs,longtable,array,enumitem,fancyhdr,titlesec}
\usepackage[colorlinks=true,linkcolor=blue!50!black,citecolor=blue!50!black,urlcolor=blue!50!black]{hyperref}
\setlength{\parindent}{1em}
\setlength{\parskip}{2pt}
\linespread{1.08}
\emergencystretch=4em
\setlist{itemsep=1pt,topsep=3pt}
\setlength{\tabcolsep}{4pt}

%%% --- status tags ---
\newcommand{\tP}{\textcolor{green!45!black}{\textbf{[Proved]}}\xspace}
\newcommand{\tC}{\textcolor{blue!55!black}{\textbf{[Lit.]}}\xspace}
\newcommand{\tV}{\textcolor{violet!80!black}{\textbf{[Verified numerically]}}\xspace}
\newcommand{\tK}{\textcolor{red!70!black}{\textbf{[Conjecture]}}\xspace}
\newcommand{\tQ}{\textcolor{orange!85!black}{\textbf{[To be verified]}}\xspace}

%%% --- theorem-like environments ---
\theoremstyle{definition}
\newtheorem{defn}{Definition}[section]
\newtheorem{thm}[defn]{Theorem}
\newtheorem{prop}[defn]{Proposition}
\newtheorem{lem}[defn]{Lemma}
\newtheorem{cor}[defn]{Corollary}
\newtheorem{ex}[defn]{Example}
\newtheorem{rem}[defn]{Remark}
\newtheorem{obs}[defn]{Observation}
\newtheorem{conj}[defn]{Conjecture}
\newcommand{\src}[1]{\par\noindent{\small\textcolor{blue!55!black}{\textbf{[Source]}~#1}}\par}

\newenvironment{hyptab}{\par\medskip\noindent{\small\textbf{Hypotheses and their roles}}\par\nopagebreak\small\noindent\begin{tabular}{>{\raggedright\arraybackslash}p{3.3cm}>{\raggedright\arraybackslash}p{5.5cm}>{\raggedright\arraybackslash}p{5.9cm}}\toprule Hypothesis & Why it is needed & Where it enters the proof\\\midrule}{\\\bottomrule\end{tabular}\par\medskip}

%%% --- notation ---
\newcommand{\R}{\mathbb R}\newcommand{\C}{\mathbb C}\newcommand{\Z}{\mathbb Z}
\newcommand{\bT}{\mathbb T}
\newcommand{\Lb}{\overline L}
\newcommand{\Jc}{\mathcal J}
\newcommand{\dbar}{\bar\partial}
\newcommand{\ch}{\mathrm{ch}}\newcommand{\Td}{\mathrm{Td}}
\newcommand{\rk}{\mathrm{rk}}
\newcommand{\HH}{\mathrm{HH}}
\newcommand{\Res}{\mathrm{Res}}
\newcommand{\dlog}{\mathrm{dlog}\,}
\newcommand{\sh}{\sharp}

\pagestyle{fancy}\fancyhf{}
\fancyhead[L]{\small Generalized holomorphic line bundles and Poisson structures}
\fancyhead[R]{\small Research note 2026-11-22}
\fancyfoot[C]{\small\thepage}
\titleformat{\section}{\large\bfseries\sffamily}{\thesection}{0.8em}{}
\titleformat{\subsection}{\normalsize\bfseries\sffamily}{\thesubsection}{0.8em}{}

\begin{document}

\begin{center}
{\LARGE\bfseries\sffamily Cohomology of generalized holomorphic line bundles on generalized complex structures from holomorphic Poisson structures\\[3pt]}
{\large Serre-type vanishing, duality and Riemann--Roch-type formulas via examples}\\[12pt]
{\large Research and writing: Claude \ ; \ Problem posing: Mori}\\[3pt]
22 November 2026 (\textbf{v2})\\[3pt]
{\small (A personal research note. It has not been peer-reviewed and is not a preprint.\\
Some proof details and bibliographic data are unverified; such places are marked explicitly by \tQ\ in the text.)}
\end{center}

\begin{abstract}
\noindent
We collect the results of an investigation of generalized-complex analogues of the ``three treasures'' of complex geometry (the Hirzebruch--Riemann--Roch theorem, Serre's vanishing theorem and Kodaira's embedding theorem), attacking the problem from the simplest non-trivial examples.
The objects are the generalized complex structure $\Jc_\sigma$ obtained from a holomorphic Poisson structure $\sigma$ by a $\beta$-transform (its type jumps along the divisor $D$ where $\sigma$ vanishes), and the generalized Dolbeault cohomology $\HH^*(L)$ of a generalized holomorphic line bundle $L$ on it.
The main results are as follows.
(1) If $\sigma$ is log-symplectic along $D$, a generalized holomorphic structure on $L$ is the same thing as a logarithmic flat connection along $D$; the residues $r_i$ are constants and $c_1(L)=-\sum r_i[D_i]$.
(2) On a surface with $D$ having only ordinary double points, $\HH^k(L)=H^k(X;\mathcal F)\oplus\C^{\mu}[-2]$, where $\mathcal F$ is a constructible complex determined by the residues and $\mu$ is the number of double points. The total Euler characteristic is $\chi(X)$, independent of $L$; the dependence on $L$ appears only in the Poincar\'e polynomial.
(3) On a projective manifold of any dimension with a $\beta$-transformed structure, for an ample Poisson module $L_0$ and $m\gg0$ the terms of transversal Dolbeault degree $a>0$ vanish and $\HH^k(L_0^m)=0$ for $k>n$ (a Serre-type theorem). Vanishing in total degree fails even for $\sigma=0$.
(4) On any compact complex manifold with any holomorphic Poisson structure, $\HH^k(L)^*\cong\HH^{2n-k}(L^\vee\otimes K_X^2)$.
(5) A weight formula in the toric case, the dependence of $\HH$ on the arithmetic of $\sigma$ in dimension 4, and computations of the local term at a triple point (partly conjectural).
A generalized Kodaira embedding theorem and a general generalized Riemann--Roch formula have not been obtained. The index-theoretic outlook (b-calculus, indicial roots) is at the level of conjecture.
\end{abstract}

\paragraph{How to read this note.}
Each statement carries one of the following tags: \tC\ quoted from the literature (source given), \tP\ proved by the author (after quoting standard facts; for the main ones a proof or proof sketch is included), \tV\ checked by computer, e.g.\ exact linear algebra (not a proof), \tK\ conjecture or outlook, \tQ\ unverified points (bibliographic data recalled from memory, or the range of applicability of a cited result not checked).
Even statements tagged \tP\ depend on the cited literature they build on. The dependencies are collected in \S\ref{sec:ledger}.

\tableofcontents

%--------------------------------------------------------------
\section{Notation}\label{sec:notation}
%--------------------------------------------------------------
Overly general symbols (wedge product, differentials, etc.) are omitted.

\begin{longtable}{>{\raggedright\arraybackslash}p{3.4cm}>{\raggedright\arraybackslash}p{11.6cm}}
\toprule Symbol & Meaning\\ \midrule\endhead
$X$ & compact complex manifold, $\dim_\C X=n$ ($n$ is even in the log-symplectic case)\\
$T,\ T^*,\ \Omega^p$ & holomorphic tangent bundle, holomorphic cotangent bundle, sheaf of holomorphic $p$-forms. $K_X=\Omega^n$\\
$\bT=T\oplus T^*$ & generalized tangent bundle. Pairing $\langle X+\xi,Y+\eta\rangle=\tfrac12(\xi(Y)+\eta(X))$, Courant (Dorfman) bracket $[X+\xi,Y+\eta]=[X,Y]+\mathcal L_X\eta-\iota_Yd\xi$\\
$\Jc$ & generalized complex structure: $\Jc^2=-1$, preserves the pairing, and its $\pm i$ eigenbundles are closed under the Courant bracket\\
$L$ & eigenbundle of $\Jc$ (maximal isotropic subbundle of $\bT_\C$, a complex Lie algebroid). $L^*\cong\Lb$\\
$\sigma$ & holomorphic Poisson structure: $\sigma\in H^0(X,\wedge^2T)$, $[\sigma,\sigma]=0$ (Schouten bracket)\\
$\sigma^\sh$ & $T^*\to T$, $\xi\mapsto\sigma(\xi,\cdot)$\\
$\Jc_\sigma$ & generalized complex structure determined by $\sigma$ ($\beta$-transform). $L_\sigma=T^{0,1}\oplus\{\sigma^\sh\xi+\xi:\xi\in T^{*1,0}\}$\\
$A=(T^*)_\sigma$ & holomorphic Lie algebroid: anchor $\sigma^\sh$, bracket $[a,b]_\sigma=\mathcal L_{\sigma^\sh a}b-\mathcal L_{\sigma^\sh b}a-d\sigma(a,b)$\\
$D=\sum D_i$ & degeneracy divisor of $\sigma$ (reduced; simple normal crossings, or only ordinary double points for surfaces)\\
$T_D=T_X(-\log D)$,\ $B=\Omega^1(\log D)$ & logarithmic tangent and cotangent bundles. $T_D$ is a Lie algebroid, $B=T_D^*$\\
$U$ & $X\setminus D$\\
log-symplectic & $\sigma\in H^0(\wedge^2T_D)$ and $\sigma^\sh:B\to T_D$ is an isomorphism\\
$L$ (as a line bundle) & to be distinguished from the eigenbundle of $\bT_\C$ by context. The line bundle is often written $E$, $L_0$, $\mathcal O(m)$, etc.\ (\S\ref{sec:setting})\\
Poisson module & holomorphic vector bundle $(E,\nabla)$ with a flat $A$-connection $\nabla$. The same thing as a generalized holomorphic (GH) bundle (Prop.~\ref{prop:ident})\\
$V$ & holomorphic vector field defined by $\nabla_\beta e=\beta(V)e$ for a local frame $e$ (Poisson vector field, Lemma~\ref{lem:flat})\\
$r_i$ & residue (a constant) of the logarithmic connection along $D_i$\\
$C(A;E)$ & $(\wedge^bT\otimes E,\ d_E=[\sigma,\cdot]+V\wedge)$ (Chevalley--Eilenberg complex of $A$)\\
$\HH^k(E)$ & $H^k(X;C(A;E))$: generalized Dolbeault cohomology (Prop.~\ref{prop:ident})\\
$E_1^{a,b}$ & $H^a(X,\wedge^bT\otimes E)$: $E_1$ term of the spectral sequence converging to $\HH^{a+b}$ ($a$ is called the transversal degree)\\
$\mathrm{DR}_{\log}(E)$ & logarithmic de Rham complex $(\Omega^\bullet(\log D)\otimes E,\nabla^{\log})$\\
$Q$ & quotient complex $C(A;E)/C(T_D\text{-subcomplex})$, supported on $D$\\
$\mu$ & (surfaces) number of ordinary double points of $D$\\
$c,\ \rho,\ b_1$ & number $c$ of irreducible components of $D$, dimension $\rho$ of the span of $[D_i]\in H^2(X;\C)$, $b_1=c-\rho$\\
$j_1:U\hookrightarrow X',\ j_2:X'\hookrightarrow X$ & open embeddings (Thm.~\ref{thm:main2})\\
$L_U$ & local system on $U$ of $\nabla$-flat sections (rank 1, monodromy $e^{2\pi ir_i}$ around $D_i$)\\
$P_E(t)$ & $\sum_k\dim\HH^k(E)\,t^k$: Poincar\'e polynomial\\
$[\,\cdot\,]$ (tables) & shorthand for $[\dim\HH^0,\dim\HH^1,\dots,\dim\HH^{2n}]$\\
\bottomrule
\end{longtable}

%--------------------------------------------------------------
\section{The classical three treasures: statements, hypotheses, and where they enter the proofs}\label{sec:classical}
%--------------------------------------------------------------
To make the target of the generalization clear, we recall the standard statements of the classical versions and the skeletons of their proofs. The content is standard textbook material and is not new.

\begin{thm}[Hirzebruch--Riemann--Roch]\tC\label{thm:HRR}
Let $M$ be a compact complex manifold and $E$ a holomorphic vector bundle. Then
$\chi(M,E)=\sum_q(-1)^q\dim H^q(M,\mathcal O(E))=\int_M\ch(E)\Td(T^{1,0}M).$
\src{Hirzebruch, \emph{Neue topologische Methoden in der algebraischen Geometrie}, Springer 1956 (projective algebraic case); Atiyah--Singer, Ann.\ Math.\ 87 (1968) 484--530, 531--545, 546--604 (general compact complex manifolds, index of the Dolbeault complex). Textbook: Huybrechts, \emph{Complex Geometry}, Springer 2005, \S5.1.}
\end{thm}
\begin{hyptab}
Compactness & The cohomology is finite-dimensional. & The Dolbeault complex is elliptic, hence Fredholm, so the index is defined.\\\midrule
Holomorphic structure & Condition $\dbar_E^2=0$ for a complex to exist. & Definition of the complex $(\Omega^{0,\bullet}(E),\dbar_E)$.\\\midrule
(Form of the conclusion) $\ch(E)$ & The index is the integral of topological data of the principal symbol. & Atiyah--Singer index theorem. $\Td$ arises from the symbol, and $E$ enters not through its rank but through $\ch(E)$.
\end{hyptab}

\begin{thm}[Serre's vanishing theorem]\tC\label{thm:Serre}
Let $X$ be a projective scheme over a field, $\mathcal O(1)$ a very ample line bundle and $\mathcal F$ a coherent sheaf. Then there is $m_0$ such that for $m\ge m_0$: (a) $H^q(X,\mathcal F(m))=0\ (q>0)$; (b) $\mathcal F(m)$ is generated by global sections.
\src{Serre, \emph{Faisceaux alg\'ebriques coh\'erents}, Ann.\ Math.\ 61 (1955) 197--278; modern proof in Hartshorne, \emph{Algebraic Geometry}, GTM 52, Springer 1977, Theorem III.5.2.}
\end{thm}
\begin{hyptab}
Coherence & Finite generation gives a surjection $\bigoplus\mathcal O(-N_i)\twoheadrightarrow\mathcal F$. & To reduce to the known computation on $\mathbb P^N$.\\\midrule
Ampleness & Twisting $\mathcal F(m)$ by large $m$ kills higher cohomology. & It is the conclusion itself.\\\midrule
Projectivity & $X$ embeds as a closed subscheme of $\mathbb P^N$, and $H^q(\mathbb P^N,\mathcal O(m))$ can be computed directly by the Cech complex. & When running the descending induction on $q$ (vanishing for $q>\dim X$) on the long exact sequence for the kernel $\mathcal K$ of the surjection.
\end{hyptab}

\begin{thm}[Kodaira--Akizuki--Nakano vanishing]\tC\label{thm:KAN}
Let $M$ be a compact K\"ahler manifold and $E$ a positive holomorphic Hermitian line bundle. If $p+q>n$ then $H^q(M,\Omega^p\otimes E)=0$. In particular $H^q(M,K_M\otimes E)=0\ (q>0)$.
\src{Kodaira, Proc.\ Nat.\ Acad.\ Sci.\ USA 39 (1953) 1268--1273; Akizuki--Nakano, Proc.\ Japan Acad.\ 30 (1954) 266--272. Textbooks: Griffiths--Harris, \emph{Principles of Algebraic Geometry}, Wiley 1978, Ch.~1 \S2; Huybrechts \S5.2.}
\end{thm}
\begin{hyptab}
Compactness & (i) Hodge theory ($H^q\cong\ker\Delta_{\dbar_E}$); (ii) no boundary terms in integration by parts. & Existence of harmonic representatives, and the integral estimate below.\\\midrule
K\"ahler & The K\"ahler identity $[\Lambda,\dbar]=-i\partial^*$. & Derivation of the Bochner--Kodaira--Nakano identity $\Delta_{\dbar_E}=\Delta_{\partial_E}+[i\Theta_E,\Lambda]$.\\\midrule
Positive ($i\Theta_E>0$) & The curvature term $[i\Theta_E,\Lambda]$ is positive definite for $p+q>n$. & For $s\in\ker\Delta_{\dbar_E}$, in $0=\|\partial_Es\|^2+\|\partial_E^*s\|^2+\langle[i\Theta_E,\Lambda]s,s\rangle$ the last term is $>0$ unless $s=0$.
\end{hyptab}

\begin{thm}[Kodaira's embedding theorem]\tC\label{thm:Kod}
Let $M$ be a compact complex manifold and $L$ a positive holomorphic line bundle. Then $M$ is projective, and for sufficiently large $m$ the map $M\to\mathbb P(H^0(M,L^m)^*)$ defined by $|L^m|$ is an embedding.
\src{Kodaira, Ann.\ Math.\ 60 (1954) 28--48; Griffiths--Harris Ch.~1 \S4; Huybrechts \S5.3.}
\end{thm}
\begin{hyptab}
$L$ positive & The curvature form $i\Theta_L$ is a K\"ahler form, so $M$ is K\"ahler. & Ensures the situation where Kodaira vanishing can be used.\\\midrule
Compactness & $H^0(M,L^m)$ is finite-dimensional, so the map to projective space is defined. & Construction of the map.\\\midrule
Enough holomorphic sections & The standard proof blows up a point $x$ to get $\tilde M$ with exceptional divisor $E_x$ and deduces surjectivity of $H^0(L^m)\to L^m\otimes\mathcal O/\mathfrak m_x^2$ from vanishing (Kodaira vanishing) of $H^1(\tilde M,\pi^*L^m\otimes\mathcal O(-2E_x))$. & Separation of points and of tangent vectors. The point is that no dimension estimate from Riemann--Roch is needed.
\end{hyptab}

\begin{rem}\tK\ (Hints for the generalization)
Classically, four items share the work: (1) Hodge theory (ellipticity), (2) K\"ahler identities $\Rightarrow$ BKN identity, (3) positivity of curvature, (4) abundance of global sections. In generalized complex geometry, (1) holds as the elliptic complex $(\wedge L^*\otimes E,d_L)$ (Prop.~\ref{prop:ellip}). The line-bundle-coefficient versions of (2) and (3) have not yet been written in the generalized K\"ahler framework. (4) does not hold in the classical form, because generalized holomorphic sections become Casimir sections of a Poisson module (\S\ref{sec:ghfun}). The present note studies these points in the simplest non-trivial examples (type-jumping $\beta$-transforms).
\end{rem}


%--------------------------------------------------------------
\section{Basics of generalized complex geometry}\label{sec:gc}
%--------------------------------------------------------------
\subsection{Definitions}
\begin{defn}[Generalized complex structure]\tC\label{def:GC}
A map $\Jc:\bT\to\bT$ with $\Jc^2=-1$ preserving the pairing is called a generalized almost complex structure. If the $+i$-eigenbundle $L\subset\bT_\C$ is closed under the Courant bracket, $\Jc$ is called integrable (a generalized complex structure).
\src{Hitchin, \emph{Generalized Calabi--Yau manifolds}, Q.\ J.\ Math.\ 54 (2003) 281--308, arXiv:math/0209099; Gualtieri, \emph{Generalized complex geometry}, Ann.\ Math.\ 174 (2011) 75--123, arXiv:math/0703298, \S3 (type Def.~3.5, canonical bundle Def.~3.7, generalized Calabi--Yau Def.~3.17, Picard group Def.~3.18; numbers checked in the source); Gualtieri's thesis arXiv:math/0401221.}
\end{defn}
\begin{ex}\tC
Complex structure $I$: $L=T^{0,1}\oplus T^{*1,0}$ (type $n$). Symplectic structure $\omega$: $L=\{X-i\omega(X)\}$ (type $0$).
\src{Hitchin 2003 (introduction); Gualtieri 2011, \S3 (examples of complex and symplectic structures; no specific example numbers could be confirmed; Def.~4.1 in \S4 is the definition of a regular point, so the ``Examples 4.1--4.2'' cited in v1 do not exist) \tQ.}
\end{ex}
\begin{defn}[Type]\tC
$\mathrm{type}(\Jc)_p=\mathrm{codim}_\C\,\pi_T(L_p)$ ($\pi_T:\bT\to T$). Points where the type is locally constant are called regular points. The type can jump in general.
\src{Gualtieri 2011, \S3, and Abouzaid--Boyarchenko, \emph{Local structure of generalized complex manifolds}, J.\ Symplectic Geom.\ 4 (2006) 43--62, arXiv:math/0511049.}
\end{defn}
\begin{prop}[Elliptic complex]\tC\label{prop:ellip}
$L$ is a complex Lie algebroid (anchor $\pi_T|_L$), and via $L^*\cong\Lb$ the exterior algebra $\wedge^\bullet L^*$ carries the Lie algebroid differential $d_L$. The complex $(\wedge^\bullet L^*,d_L)$ is elliptic, and on a compact manifold the generalized Dolbeault cohomology $H^k(\wedge L^*,d_L)$ is finite-dimensional. Using the spinor bundle $K$, under $\wedge^{n-k}\Lb\cdot K\cong U^k$ the differential $d_L$ corresponds to $\dbar:U^k\to U^{k-1}$.
\src{Gualtieri 2011, Prop.~3.13 ($(\wedge L^*,d_L)$ is elliptic), Cor.~3.14 (finite-dimensional cohomology when compact). In the proof in \S\ref{sec:duality}, for $\Jc_\sigma$ ellipticity is verified directly from the principal symbol (\tP).}
\end{prop}

\subsection{The $\beta$-transform $\Jc_\sigma$ and generalized holomorphic bundles}\label{sec:setting}
Let $X$ be a complex manifold and $\sigma$ a holomorphic Poisson structure.

\begin{prop}\tC\label{prop:Jsigma}
(a) $L_\sigma=T^{0,1}\oplus\{\sigma^\sh\xi+\xi:\xi\in T^{*1,0}\}$ is a maximal isotropic subbundle of $\bT_\C$ and is integrable. The corresponding generalized complex structure $\Jc_\sigma$ has type $n-\mathrm{rk}\,\sigma^\sh_p$ at $p\in X$ ($\mathrm{rk}$ is the complex rank). In particular it has type 0 (symplectic type) where $\sigma$ is non-degenerate, and type $n$ (complex type) where $\sigma$ vanishes.
(b) $(\Xi,\xi)\mapsto(\Xi+\sigma^\sh\xi,\xi)$ is an isomorphism of complex Lie algebroids $T^{0,1}\bowtie(T^{1,0})^*_\sigma\to L_\sigma$ (the left side is a matched pair, written $T^{0,1}\bowtie A^{1,0}$ in LGSX).
\src{(a) Gualtieri, \emph{Branes on Poisson varieties}, arXiv:0710.2719, Prop.~5; Laurent-Gengoux--Sti\'enon--Xu (LGSX below), \emph{Holomorphic Poisson manifolds and holomorphic Lie algebroids}, IMRN 2008, Art.\ ID rnn088, arXiv:0707.4253, Thm.~2.7 and Rem.~2.8 (there written $\Jc_\pi$, $\pi=\pi_R+i\pi_I$). The type formula is the author's reorganization, consistent with Gualtieri's examples ($\mathbb{CP}^2$ etc.) (\tQ). (b) LGSX Thm.~4.21 and Lemma~4.22 (the factor 4 in $\Jc_{4\pi}$ comes from the normalization of $\pi=\pi_R+i\pi_I$, LGSX Prop.~3.7). Gualtieri's paper and LGSX differ in the choice of sign of eigenvalues ($\ker(\Jc\mp i)$), but the two are exchanged by complex conjugation, so the statements below are unaffected.}
\end{prop}

\begin{defn}[Generalized holomorphic bundle]\tC
A complex vector bundle $E$ equipped with a flat $L$-connection $\dbar_E:\Gamma(E)\to\Gamma(L^*\otimes E)$ ($\dbar_E^2=0$).
\src{Gualtieri 2011, \S3 (generalized holomorphic bundles are introduced in prose after Thm.~3.15, with no numbered definition; the complex case is Example 3.19; Def.~3.18 is the definition of the Picard group); Gualtieri, \emph{Branes on Poisson varieties}, \S4.1.}
\end{defn}
\begin{defn}[Poisson module]\tC
A holomorphic vector bundle $E$ with a flat $A$-connection $\nabla:\Gamma(E)\to\Gamma(A^*\otimes E)$, i.e.\ satisfying $\nabla_{\xi}(fs)=f\nabla_\xi s+(\sigma^\sh\xi)(f)s$ and $[\nabla_a,\nabla_b]=\nabla_{[a,b]_\sigma}$ ($a,b$ holomorphic 1-forms). Here $A=(T^*)_\sigma$.
\src{LGSX, Def.~4.15 and Lemma~4.16; Gualtieri, \emph{Branes on Poisson varieties}, \S2--4.}
\end{defn}

\begin{prop}[Generalized holomorphic bundles $=$ Poisson modules, and the Dolbeault model]\label{prop:ident}
Let $(X,\sigma)$ be as above.
\begin{enumerate}
\item[(a)] A generalized holomorphic bundle with respect to $\Jc_\sigma$ is the same thing as a holomorphic vector bundle with a Poisson module structure. \tC\ Gualtieri, \emph{Branes on Poisson varieties}, Prop.~8; LGSX Lemma~4.16. A direct proof is given below (\tP). For line bundles, by Prop.~9 of the same paper, the obstruction is a class $\sigma\cdot\alpha$ in $H^1(T^{1,0})$ together with a class in $H^2_\sigma$, and when structures exist they form an affine space over $H^1_\sigma(X)=\HH^1(\mathcal O)$. \tC
\item[(b)] Since $L^*=T^{*0,1}\oplus T^{1,0}$, we have $\Gamma(\wedge^kL^*\otimes E)=\bigoplus_{a+b=k}A^{0,a}(\wedge^bT\otimes E)$ and $d_L=\dbar+(-1)^ad_A$, where $d_A=[\sigma,\cdot]+V\wedge$ acts on the coefficients in a holomorphic frame. Hence
\[H^k_{GC}(E)=H^k\bigl(\Gamma(\wedge^\bullet L^*\otimes E),d_L\bigr)=\mathbb H^k\bigl(X;(\wedge^\bullet T\otimes E,\ [\sigma,\cdot]+V\wedge)\bigr)=:\HH^k(E).\]
\tC\ This is LGSX Thm.~4.19 ($H^*(A,E)=H^*(T^{0,1}\bowtie A^{1,0},E)$, with $H^*(A,E)$ defined as sheaf cohomology in Def.~4.17). \tP\ The structure of the proof is given below, but it depends on Mokri's theorem on matched pairs (Matched pairs of Lie algebroids, Glasgow Math.\ J.\ 39 (1997) 167--181, doi:10.1017/S0017089500032055; year and DOI checked, volume/pages from memory \tQ; LGSX Lemma~4.16: a complex vector bundle is an $A$-module iff it is a $T^{0,1}X\bowtie A^{1,0}$-module; the proof of Thm.~4.19 is only stated to be similar to that of Thm.~4.11) and on LGSX Lemma~4.22 (not re-derived).
\end{enumerate}
\end{prop}
\begin{proof}[Proof of (a) (sketch)]
($\Leftarrow$) Given a Poisson module structure and a holomorphic frame $\{e_s\}$, put $\nabla_{\bar\Xi}\sum f_se_s=\sum(\bar\Xi f_s)e_s$ (holomorphic structure) and $\nabla_\xi\sum f_se_s=\sum((\sigma^\sh\xi)f_s)e_s+f_s\nabla_\xi e_s$ ($\xi\in A^{1,0}$). This is the unique extension forced by the Leibniz rule, and it is well defined because $\nabla_\xi e_s$ is holomorphic. Flatness: the $(\bar\Xi,\bar\Upsilon)$ components follow from integrability of the holomorphic structure; the $(\bar\Xi,\xi)$ components are equivalent to ``$\nabla_\xi$ maps holomorphic sections to holomorphic sections'', which holds by construction; the $(\xi,\eta)$ components are the $C^\infty$-linear extension of flatness of the Poisson module (for holomorphic $\xi,\eta$) (the curvature is a tensor).
($\Rightarrow$) Restricting a flat $L$-connection to $T^{0,1}$ gives a holomorphic structure (Koszul--Malgrange); flatness of mixed type shows that $\nabla_\xi$ maps holomorphic sections to holomorphic sections, and the module conditions follow from the Leibniz rule and flatness.\qedhere
\end{proof}
\begin{proof}[Structure of the proof of (b)]
By the general theory of matched pairs (Mokri), the Chevalley--Eilenberg complex of $L\cong T^{0,1}\bowtie A^{1,0}$ (Prop.~\ref{prop:Jsigma}(b)) is the total complex of $C(T^{0,1};\,C(A^{1,0};E))$. Here the $T^{0,1}$ direction is the Dolbeault resolution of $\wedge^bT\otimes E$ (differential $\dbar$), and the $A^{1,0}$ direction is $C(A;E)=(\wedge^bT\otimes E,\ [\sigma,\cdot]+V\wedge)$. The two differentials anti-commute with the Koszul sign $(-1)^a$. The total complex is a fine resolution of the complex $(\wedge^\bullet T\otimes E,d_A)$, so its cohomology is the hypercohomology.\qedhere
\end{proof}

\begin{rem}[\tC\ Note on integrability]
A notion of Lie groupoid integrating the complex Lie algebroid $L=L_\sigma$ itself was not found in the literature surveyed. Thms.~3.17 and 3.22 of LGSX arXiv:0803.2031 relate integrability of a holomorphic Lie algebroid $A$ to that of the real Poisson structures $\mathrm{Re}\,\sigma$, $\mathrm{Im}\,\sigma$, but do not treat $T^{0,1}\bowtie A^{1,0}$ (Rem.~4.14 of that paper only poses it as a question). Thm.~7.1 of Bailey--Gualtieri, arXiv:1611.03850 (title \tQ), states the equivalence between a generalized complex structure integrating to a weak holomorphic symplectic groupoid (a Morita-class notion) and integrability of its underlying real Poisson structure. \tP\ For the smooth cubic $\sigma$ on $\mathbb{CP}^2$, the leaves of $\mathrm{Im}\,\sigma$ are the open leaf $U$ (isotropy $0$) and the points of $D$ (isotropy $\mathrm{aff}(1,\C)$, with trivial centre); by Crainic--Fernandes, Ann.\ Math.\ 157 (2003) 575--620 (arXiv:math/0105033) Cor.~5.2 (``if all $N_x$ are trivial, then integrable'') $\mathrm{Im}\,\sigma$ is integrable. The hypotheses of Bailey--Gualtieri were checked only through the summary obtained \tQ.
\end{rem}

\subsection{On generalized holomorphic sections}\label{sec:ghfun}
\begin{obs}\tP\label{obs:casimir}
A generalized holomorphic section ($\dbar_Es=0$) of a Poisson module $(E,\nabla)$ is a flat (Casimir-type) section: $\HH^0(E)=\{s\in H^0(X,E):\nabla_\xi s=0\ (\forall\xi)\}$. Near a point of type $0$ ($\sigma$ non-degenerate), $\nabla$ corresponds via $\sigma^\sh$ to an ordinary flat connection, and flat sections are locally constant.
\end{obs}
\begin{proof}
$\HH^0=\ker(d_A:H^0(E)\to H^0(T\otimes E))$, and $d_As=(\xi\mapsto\nabla_\xi s)$ (Prop.~\ref{prop:ident}(b)).
\end{proof}
\begin{rem}\tK
This means that, in the regular case of non-zero type as well (the Pal--Poddar framework), generalized holomorphic sections are locally constant in the leaf directions, so they cannot give an embedding of $M$ into projective space. A generalized Kodaira embedding theorem will have to be reformulated as a map to the leaf space (or to the quotient separated by Casimir sections). The present note does not treat the generalization of the Kodaira embedding theorem itself.
\end{rem}


%--------------------------------------------------------------
\section{The log-symplectic case: Poisson modules $=$ logarithmic flat connections}\label{sec:logsymp}
%--------------------------------------------------------------
\subsection{Setting}
Let $X$ be a compact complex manifold ($\dim_\C X=n=2k$), $\sigma$ a holomorphic Poisson structure, and $D\subset X$ a reduced divisor (simple normal crossings; for surfaces, ordinary double points are also allowed).
Put $T_D=T_X(-\log D)$, $B=\Omega^1(\log D)$.

\begin{defn}[Log-symplectic]\tC
$\sigma$ is called log-symplectic along $D$ if $\sigma\in H^0(\wedge^2T_D)$ and $\sigma^\sh:B\to T_D$ is an isomorphism (equivalently, $\omega:=\sigma^{-1}\in H^0(\Omega^2(\log D))$ is non-degenerate). Then $\mathrm{div}(\sigma^k)=D$, i.e.\ $-K_X=\mathcal O(D)$.
\src{Goto, \emph{Unobstructed deformations of generalized complex structures induced by $C^\infty$ logarithmic symplectic structures}, arXiv:1501.03398 (journal and pages unverified \tQ); Gualtieri--Li, \emph{Symplectic groupoids of log symplectic manifolds}, IMRN 2014, arXiv:1206.3674; Pym, \emph{Constructions and classifications of projective Poisson varieties}, Lett.\ Math.\ Phys.\ 108 (2018) 573--632 (checked). The journal of Gualtieri--Li and the publication data of Goto are unverified \tQ.}
For surfaces this means that $\sigma\in H^0(-K_X)$ has a reduced zero divisor with only ordinary double points. Locally $\sigma=h\,\partial_1\wedge\partial_2$ with $h=zu$ (smooth points) or $h=z_1z_2u$ (double points), $u$ a unit.
\end{defn}

\subsection{Statement}
\begin{thm}[Poisson modules $=$ logarithmic flat connections]\tP\label{thm:main1}
Let $(X,D,\sigma)$ be log-symplectic (any even dimension, $D$ simple normal crossings) and $E$ a holomorphic line bundle carrying a Poisson module structure.
\begin{enumerate}
\item[(a)] For a local frame $e$, a holomorphic vector field $V$ is defined by $\nabla_\beta e=\beta(V)e$, and $[\sigma,V]=0$ (a Poisson vector field). Under a change of frame $e\mapsto ge$, $V$ changes only by a Hamiltonian term of the form $X_{\log g}$.
\item[(b)] $V$ is automatically logarithmic: $V\in T_D$. Hence the $A$-module structure extends uniquely to a $B$-module structure, and via the isomorphism $\sigma^\sh:B\cong T_D$ it is the same thing as a flat connection $\nabla^{\log}:E\to\Omega^1(\log D)\otimes E$ with logarithmic poles.
\item[(c)] The residue of $\nabla^{\log}$ along $D_i$ is a constant $r_i\in\C$, and $c_1(E)=-\sum_ir_i[D_i]$ in $H^2(X;\C)$.
\item[(d)] Near a point of $D$ one can choose a frame with $\nabla^{\log}e=\sum_ir_i\,\dlog z_i\,e$.
\item[(e)] The subcomplex $C(T_D;E)=(\wedge^bT_D\otimes E,d_E)\subset C(A;E)$ is isomorphic via $\wedge^b\sigma^\sh$ to the logarithmic de Rham complex $\mathrm{DR}_{\log}(E)$.
\end{enumerate}
\end{thm}
\begin{hyptab}
Log-symplectic ($\sigma^\sh:B\cong T_D$) & To identify the Lie algebroid structure $(T_D^*)_\sigma$ on $B$ with $T_D$, and to treat Poisson modules as ordinary flat connections (with logarithmic poles). & Lemma~\ref{lem:iso}, Theorem (b), (e).\\\midrule
$D$ reduced, simple normal crossings & To estimate the pole order from each factor of the local equation $f=z_1\cdots z_s$ independently. & Lemma~\ref{lem:log} (deducing $V(z_i)\in(z_i)$ from $V(f)\in(f)$).\\\midrule
$X$ compact & For the relation between $c_1(E)$ and the residues (an equality of global cohomology classes). & (c). That the residues are constant is local.\\\midrule
$\sigma$ Poisson ($[\sigma,\sigma]=0$) & $A$ becomes a Lie algebroid and $d_E^2=0$. & Lemma~\ref{lem:flat}.
\end{hyptab}

\begin{lem}[Flatness $\Leftrightarrow$ Poisson vector field]\tP\label{lem:flat}
For a local frame $e$ put $\Theta(\beta):=\nabla_\beta e/e$. Then $\nabla_a\nabla_be=(\sigma^\sh a)(\Theta b)e+\Theta(a)\Theta(b)e$. Hence flatness $[\nabla_a,\nabla_b]=\nabla_{[a,b]_\sigma}$ is equivalent to
$(\sigma^\sh a)(\Theta b)-(\sigma^\sh b)(\Theta a)=\Theta([a,b]_\sigma)$, i.e.\ $\Theta$ is a 1-cocycle of the Lichnerowicz complex: $\Theta(\beta)=\beta(V)$, $[\sigma,V]=0$.
In local coordinates with $\sigma=h\partial_1\wedge\partial_2$, $a=\Theta(dz_1)$, $b=\Theta(dz_2)$, the condition reads $h(a_1+b_2)=h_1a+h_2b$, i.e.\ $V(h)=h\,\mathrm{div}V$. The differential $d_E$ is, on $C^0$, $d_E(fe)=(-\sigma^\sh(df)+fV)e$, and for $c=X\otimes e$, $X=p\partial_1+q\partial_2$,
\[(d_Ec)(dz_1,dz_2)=h(p_1+q_2)-h_1p-h_2q+qa-pb.\]
\end{lem}
\begin{lem}[$V$ is logarithmic]\tP\label{lem:log}
Let $f=z_1\cdots z_s$ be a local equation of $D$. From $L_V\sigma=0$ we get $L_V\sigma^k=0$. Since $\sigma^k=f'\partial_1\wedge\cdots\wedge\partial_n$ ($f'=f\cdot$unit), $V(f')-f'\,\mathrm{div}V=0$, hence $V(f)\in(f)$. Now $V(f)/f=\sum_iV(z_i)/z_i$, and the $i$-th term has a pole only along $\{z_i=0\}$. Since the whole is holomorphic, each $V(z_i)/z_i$ is holomorphic, i.e.\ $V(z_i)\in(z_i)$.
\end{lem}
\begin{lem}[$\sigma^\sh:B\to T_D$ is an isomorphism of Lie algebroids]\tP\ (standard)\label{lem:iso}
$\sigma\in\wedge^2T_D$ is a non-degenerate Poisson structure on the Lie algebroid $T_D$. That $\omega=\sigma^{-1}$ is $d_{T_D}$-closed follows from $\omega$ being closed on $U$ (the inverse of a Poisson structure) and from the fact that a section of $\Omega^3(\log D)$ vanishing on a dense open set is $0$. For a non-degenerate Poisson structure coming from a closed 2-form, $\sigma^\sh$ preserves brackets. Hence $(T_D^*)_\sigma=B\cong T_D$.
\end{lem}
\begin{lem}[Residues, local normal form, $c_1$]\tP\ /\ \tC\label{lem:res}
Let $\nabla^{\log}e=\theta e$ with $\theta\in\Omega^1(\log D)$; $\theta$ is closed by flatness. The residue $\theta_i|_{D_i}$ along $D_i$ is seen to be constant by restricting $d\theta$ to $D_i$. Since $\theta-\sum r_i\dlog z_i$ is a closed holomorphic 1-form, it equals $dg$, and for the frame $e'=e^{-g}e$ we get $\nabla^{\log}e'=\sum r_i\dlog z_i\,e'$ (this is (d)). The relation $c_1(E)=-\sum r_i[D_i]$ is the Esnault--Viehweg residue formula (\tC\ Esnault--Viehweg, \emph{Logarithmic de Rham complexes and vanishing theorems}, Invent.\ Math.\ 86 (1986) 161--194, \S2; Deligne, \emph{\'Equations diff\'erentielles \`a points singuliers r\'eguliers}, LNM 163, Springer 1970, Ch.~II. Sign check: for $\mathcal O(D)$ with the canonical section flat, the residue is $-1$ and $c_1=[D]$).
The set of all module structures with the same $c_1$ is a torsor under $H^0(\Omega^1(\log D))=\ker(\C^c\to H^2(X;\C))$ (Deligne), of dimension $c-\rho$. If $D$ is irreducible the structure is unique; if reducible, it forms an affine family of residue vectors.
\end{lem}
\begin{proof}[Proof of Theorem~\ref{thm:main1}]
(a) Lemma~\ref{lem:flat}. (b) By Lemma~\ref{lem:log}, $\Theta$ extends to $B$, and since flatness is an equality on a dense open set it persists after the extension. Combine with Lemma~\ref{lem:iso}. (c), (d) Lemma~\ref{lem:res}. (e) Lemma~\ref{lem:iso}.
\end{proof}

%--------------------------------------------------------------
\section{The simplest non-trivial example: $\mathbb{CP}^2$ and a smooth cubic}\label{sec:cp2}
%--------------------------------------------------------------
Let $F$ be a smooth homogeneous cubic polynomial on $\mathbb C^3$ and $\Pi^{ij}=\epsilon_{ijk}\partial_kF$. Then $\Pi$ is a homogeneous Poisson structure (of weight 0) whose contraction with the Euler vector field is $0$, and it defines a holomorphic Poisson structure $\sigma\in H^0(\wedge^2T)=H^0(\mathcal O(3))$ on $\mathbb{CP}^2$. $\sigma$ vanishes to first order along $D=\{F=0\}$ (a smooth cubic curve). This is a standard construction (identification of the source \tQ; the vanishing order along $D$ was checked by a local computation). Below $E=\mathcal O(m)$.

\subsection{Which $\mathcal O(m)$ are generalized holomorphic}
\begin{prop}\tP\label{prop:cp2mod}
For every $m\in\Z$ the bundle $\mathcal O(m)$ has exactly one Poisson module structure, and the residue of its logarithmic connection is $r=-m/3$.
\end{prop}
\begin{proof}
By Theorem~\ref{thm:main1}(c), $r=-m/3$ (as $[D]=3H$). Existence: vanishing of the Atiyah class of $\mathcal O(m)$ in $H^1(\Omega^1(\log D))$ is equivalent to $c_1\in\C[D]$ (exact sequence $H^0(\mathcal O_D)\to H^1(\Omega^1)\to H^1(\Omega^1(\log D))$), and $m[H]\in\C\cdot3[H]$. The curvature is an element of $H^0(\Omega^2(\log D))=\C\omega$ with residue $dr=0$, hence zero (flat). Uniqueness: the difference of two logarithmic connections lies in $H^0(\Omega^1(\log D))=0$ (the connecting map sends $1\mapsto[D]=3H\ne0$).\qedhere
\end{proof}
Another view: for the Casimir $F$ of $\Pi$, $F^{m/3}$ gives a (multivalued) flat section of $\mathcal O(m)$ on $U=\mathbb{CP}^2\setminus D$. $\mathcal O(m)|_U$ is a rank-one local system $L_m$ with monodromy $e^{-2\pi im/3}$ around $D$. The period 3 corresponds to $\pi_1(U)=\Z/3$. \tP\ Locally $\nabla e=-\tfrac m3\dlog z\,e$.

\subsection{Result}
\begin{thm}\tP\label{thm:cp2}
The dimensions $[\dim\HH^0,\dots,\dim\HH^4]$ of the generalized Dolbeault cohomology
\[\HH^*(\mathcal O(m))=\mathbb H^*\bigl(\mathbb{CP}^2;(\wedge^\bullet T(m),d_\sigma)\bigr)\]
are as follows.
\begin{center}\small
\begin{tabular}{lll}\toprule
$m$ & dimensions & Interpretation\\\midrule
$m\in3\Z_{\ge0}$ & $[1,0,2,0,0]$ & $H^*(U;\C)$. $\HH^0=\C F^{m/3}$ is a Casimir section\\
$m\in3\Z_{\le-6}$ & $[0,0,2,0,1]$ & $H^*_c(U;\C)$\\
otherwise (including $m=-3$) & $[0,0,3,0,0]$ & $H^*(U;L_m)$ ($m\notin3\Z$); $m=-3$ is special\\\bottomrule
\end{tabular}
\end{center}
The total Euler characteristic is always $3=\chi(\mathbb{CP}^2)$.
\end{thm}
\begin{hyptab}
$D$ smooth, $\sigma$ vanishing to first order along $D$ & The differential of the quotient complex $Q$ is a constant multiple of the isomorphism $d\sigma|_D:N_D\to\wedge^2T|_D$. & Step (3) of the proof. $Q$ is acyclic for $m\ne-3$.\\\midrule
$\pi_1(U)=\Z/3$ (Milnor fibre $\simeq$ a bouquet of 8 copies of $S^2$) & To determine the cohomology of $U$ explicitly. & Right-hand side of the table in step (2).\\\midrule
$m\ne-3$ & The differential $-(1+m/3)\cdot(d\sigma|_D)$ of $Q$ is not the zero map. & Only $m=-3$ is treated separately (step (4)).
\end{hyptab}
\begin{proof}[Steps of the proof]
(1) By Proposition~\ref{prop:cp2mod} and Theorem~\ref{thm:main1}(e), the subcomplex $C_{\log}=(\wedge^bT_D\otimes\mathcal O(m))\cong\mathrm{DR}_{\log}$.
(2) In a local frame $\nabla^{\log}e=r\,\dlog z\,e$ ($r=-m/3$). At the stalk at a point of $D$, the logarithmic de Rham complex acts on the monomial $z^ne$ by the factor $(n+r)=(n-m/3)$. The stalk cohomology is $[1,1]$ for $m\in3\Z_{\ge0}$ and $0$ otherwise. Hence for $m\in3\Z_{\ge0}$ the natural map $\mathrm{DR}_{\log}\to Rj_*L_m$ is a stalkwise quasi-isomorphism ($\mathbb H(C_{\log})=H^*(U;\C)$), and otherwise $j_!L_m\to\mathrm{DR}_{\log}$ is a quasi-isomorphism ($\mathbb H(C_{\log})=H_c^*(U;L_m)$). Topology of $U$: the Milnor fibre $F=1$ is a bouquet of $\mu=8$ copies of $S^2$, $U=\{F=1\}/\mu_3$, and $H^2=\C^8=2+3+3$ (trivial character and two non-trivial characters). Since $U$ is affine, $H^k=0$ for $k>2$. Hence $H^*(U)=[1,0,2,0,0]$, $H^*(U;L)=[0,0,3,0,0]$ for a non-trivial character, and by Verdier duality $H^*_c(U;\C)=[0,0,2,0,1]$.
(3) The quotient $Q=C(A;\mathcal O(m))/C_{\log}$ is a complex on $D$ with $Q^1=N_D(m)=\mathcal O_D(m+3)$, $Q^2=\wedge^2T(m)|_D=\mathcal O_D(m+3)$. Setting $z=0$ in the formula of Lemma~\ref{lem:flat} gives $d_Q=-u(1-r)p$, and in invariant form $d_Q=-(1+m/3)\cdot(d\sigma|_D)$. Since $d\sigma|_D$ is an isomorphism, $Q$ is acyclic for $m\ne-3$, and $\HH^*(\mathcal O(m))=\mathbb H^*(C_{\log})$.
(4) $m=-3$: the terms $E_1^{a,b}=H^a(\wedge^bT(-3))$ are 1-dimensional at $(2,0),(1,1),(0,2)$ (Bott); all lie in total degree 2, so all further differentials vanish and the result is $[0,0,3,0,0]$. No sign check is needed.
\end{proof}
\begin{rem}[\tV]
The table above agrees with a direct computation, for $m=-15,\dots,15$ and the Fermat and Hesse cubics, of the differential $d_1$ of the spectral sequence $E_1^{a,b}=H^a(\wedge^bT(m))$ by exact linear algebra modulo the prime $32003$ (for $m\ge-2$ as a complex of global sections, for $m\le-4$ as a Cech computation with Laurent monomials) (scripts: \texttt{cp2/poisson.py, table.py, neg\_check.py, predict.py}).
\end{rem}
\begin{rem}[\tC\ Related work]
Eq.~(11), $0\to T^*\xrightarrow{\sigma}T\to T|_D\to0$, and Prop.~8 of Hitchin, \emph{Deformations of holomorphic Poisson manifolds}, arXiv:1105.4775 (Moscow Math.\ J.\ 12 (2012) 567--591) are related, but this result for $\mathbb{CP}^2$ itself is stated neither in that paper nor in Gualtieri's (as far as surveyed).
\end{rem}

\begin{obs}[How the dependence on $m$ appears]\tP
The total Euler characteristic does not depend on $m$ (a topological invariant). The dependence on $m$ appears only in the Poincar\'e polynomial, in particular in the period-3 jumps and in $\HH^0$ (the Casimir section $F^{m/3}$). Even for $m\gg0$, $\HH^2$ stays 2- or 3-dimensional, so Serre vanishing in total degree does not hold.
\end{obs}


%--------------------------------------------------------------
\section{Main theorem (surfaces): generalized Dolbeault cohomology of log-symplectic surfaces}\label{sec:main2}
%--------------------------------------------------------------
\subsection{Statement}
Let $X$ be a compact complex surface, $\sigma\in H^0(-K_X)$ with reduced zero divisor $D=\bigcup D_i$ (irreducible components $D_i$, with only ordinary double points as singularities of $D$, their number being $\mu$), and $E$ a line bundle with a Poisson module structure. The residues $r_i$ (constant along each branch) are given by Theorem~\ref{thm:main1}(c). (The branches at a double point are regarded as different branches.)
Let $I_*=\{i:r_i\in\Z_{\le0}\}$, $I_!=\{i:r_i\notin\Z_{\le0}\}$, $X'=X\setminus\bigcup_{i\in I_!}D_i$, $j_1:U\hookrightarrow X'$, $j_2:X'\hookrightarrow X$, and put
\[\mathcal F:=j_{2!}\,Rj_{1*}L_U\]
($Rj_*$ along the components in $I_*$, extension by zero along those in $I_!$).

\begin{thm}[Main theorem for surfaces]\tP\label{thm:main2}
In the above setting assume all $r_i\neq1$. Then
\[\HH^k(E)\;\cong\;H^k(X;\mathcal F)\ \oplus\ \C^{\mu}\ (\text{only for }k=2).\]
In particular
\begin{enumerate}
\item[(i)] if all $r_i\notin\Z_{>0}$ then $\HH^*(E)=H^*(U;L_U)\oplus\C^\mu[-2]$;
\item[(ii)] if all $r_i\notin\Z_{\le0}$ then $\HH^*(E)=H^*_c(U;L_U)\oplus\C^\mu[-2]$;
\item[(iii)] the total Euler characteristic is $\sum(-1)^k\dim\HH^k(E)=\chi(U)+\mu=\chi(X)$ (independent of $E$).
\end{enumerate}
\end{thm}
\begin{hyptab}
Surface, $\sigma\in H^0(-K_X)$, $D$ reduced with only ordinary double points & To use the local normal forms $h=zu$ or $z_1z_2u$ and compute the quotient complex $Q$ explicitly. & Lemma~\ref{lem:Q}, Prop.~\ref{prop:moser}, Lemma~\ref{lem:conn}.\\\midrule
$r_i\ne1$ (all branches) & If $r_i=1$ the differential $-u(1-r)p$ of $Q$ vanishes, the cohomology sheaves of $Q$ become infinite-dimensional, and the reduction breaks down. & Lemma~\ref{lem:Q}. ($r_i=1$ is the special case $E=K_X$, Prop.~\ref{prop:K}.)\\\midrule
Residues not positive integers / not non-positive integers & Condition for the comparison of the logarithmic complex with $Rj_*L_U$, $j_!L_U$ (Deligne) to be a quasi-isomorphism. & Lemma~\ref{lem:stalk}.\\\midrule
$X$ compact & $R\Gamma(X;j_!L_U)=R\Gamma_c(U;L_U)$ and the Euler characteristic computation. & (ii), (iii).\\\midrule
$-K_X=\mathcal O(D)$ (anticanonical curve) & Gives $\chi(D)=\mu$ ($\omega_D=\mathcal O_D$, $p_a(D)=1$). & (iii).
\end{hyptab}

\subsection{Proof}
\begin{lem}[Stalks of the logarithmic complex and Deligne's comparison]\tP\ /\ \tC\label{lem:stalk}
In a local frame (Theorem~\ref{thm:main1}(d)), the stalk of $\mathrm{DR}_{\log}$ at a point of $D_I$ ($|I|=s$) is the tensor product of the Koszul complex of the commuting operators $z_i\partial_i+r_i\ (i\in I)$ with the ordinary de Rham complex in the remaining variables. On a monomial $z^a$, $z_i\partial_i+r_i=a_i+r_i$ (a scalar). Hence the stalk cohomology is $\wedge^*\C^s$ (dimension $2^s$, Euler characteristic 0) exactly when $-r_i\in\mathbb N$ (i.e.\ $r_i\in\Z_{\le0}$) for all $i\in I$, and $0$ otherwise. On the other hand the stalk of $Rj_*L_U$ at the same point is $\wedge^*\C^s$ exactly when all $e^{2\pi ir_i}=1$ ($r_i\in\Z$). Therefore:
\begin{itemize}
\item the natural map $\mathrm{DR}_{\log}\to Rj_*L_U$ is a quasi-isomorphism when no $r_i$ is a positive integer (Deligne, LNM 163, Ch.~II \S3, number \tQ; here confirmed by direct computation);
\item the extension by zero $j_!L_U\to\mathrm{DR}_{\log}$ (a zero extension of flat sections is holomorphic) is a quasi-isomorphism when no $r_i$ lies in $\Z_{\le0}$;
\item in the mixed case, over $X'$ we have $\mathrm{DR}_{\log}\simeq Rj_{1*}L_U$, and over $\bigcup_{I_!}D_i$ the stalks are $0$, so $\mathrm{DR}_{\log}\simeq\mathcal F$.
\end{itemize}
Division by the coefficients $1/(a_i+r_i)$ does not destroy convergence, since $|a_i+r_i|\ge c>0$ except for finitely many $a$.
\end{lem}
\begin{lem}[The quotient complex $Q$]\tP\ \tV\label{lem:Q}
$Q:=C(A;E)/C(T_D;E)$ is supported on $D$ and lives only in degrees 1, 2: $Q^1=(T/T_D)\otimes E$, $Q^2=(\wedge^2T/\wedge^2T_D)\otimes E$.
\begin{itemize}
\item Smooth points ($h=zu$): $Q^1=\partial_z\otimes\mathcal O_D$, $Q^2=\partial_z\wedge\partial_w\otimes\mathcal O_D$, $d_Q(p)=-u(1-r)\,p|_D$. Hence acyclic when $r\ne1$.
\item Double points ($h=z_1z_2u$): $Q^1=\mathcal O_{D_1}\partial_1\oplus\mathcal O_{D_2}\partial_2$, $Q^2=\mathcal O_D$, $d_Q(p,q)=-u[(1-r_1)z_2p+(1-r_2)z_1q]\bmod z_1z_2$. If $r_1,r_2\ne1$ it is injective with image the maximal ideal $\mathfrak m_D$, so $H^1(Q)_p=0$ and $H^2(Q)_p=\mathcal O_D/\mathfrak m_D=\C$.
\end{itemize}
Hence for $r_i\ne1$, $Q\simeq\bigoplus_{\text{double points }p}\C_p[-2]$ in the derived category.
Derivation: restrict the formula of Lemma~\ref{lem:flat} to $z=0$ and substitute $a|_D=\theta_2h$, $b|_D=-\theta_1u$, $\theta_1|_D=r$ (Theorem~\ref{thm:main1}(d)). (\tV\ With random polynomials $u,g,p,q$ and residues, at smooth points and double points (20 trials each), $d_1d_0=0$ and the formula for $d_Q$ were confirmed in exact arithmetic. Script \texttt{cp2/ncd\_cochain\_check.py}.)
\end{lem}
\begin{prop}[Logarithmic Moser normal form at a double point]\tP\label{prop:moser}
Near a double point $p$ let $\omega=\sigma^{-1}=u^{-1}\dlog z_1\wedge\dlog z_2$ and $c:=u(0)$. There is a biholomorphic map $\phi$ preserving each branch of $D$ and fixing $p$ such that $\phi^*\omega=c^{-1}\dlog z_1\wedge\dlog z_2$. The constant $c$ is a local invariant and cannot in general be made $1$.
\end{prop}
\begin{proof}
Let $\omega_0=\dlog z_1\wedge\dlog z_2$, $k=1/u-1/c$ ($k(0)=0$), $\omega_t=(c^{-1}+tk)\omega_0$. $\omega_t$ is closed and non-degenerate near $p$ for each $t\in[0,1]$. For $k=\sum k_{ab}z_1^az_2^b$,
\[\alpha=\sum_{a\ge1,b\ge0}\tfrac{k_{ab}}{a}z_1^az_2^b\,\dlog z_2-\sum_{b\ge1}\tfrac{k_{0b}}{b}z_2^b\,\dlog z_1\]
is a convergent logarithmic 1-form with $d\alpha=k\omega_0$ (using $d(z^{ab}\dlog z_2)=az^{ab}\omega_0$, $d(z_2^b\dlog z_1)=-bz_2^b\omega_0$). Since $X\mapsto\iota_X\omega_t$ is an isomorphism $T_D\to\Omega^1(\log D)$, there is a unique $X_t\in T_D$ with $\iota_{X_t}\omega_t=-\alpha$; it is tangent to each branch and vanishes at $p$. Its flow $\phi_t$ preserves $D$, and $\tfrac{d}{dt}\phi_t^*\omega_t=\phi_t^*(d\iota_{X_t}\omega_t+k\omega_0)=\phi_t^*(-d\alpha+k\omega_0)=0$. Hence $\phi_1^*\omega_1$ is of the form $\omega_{c^{-1}}$.\qedhere
\end{proof}
\begin{lem}[Vanishing of the connecting morphism]\tP\ \tV\label{lem:conn}
For the exact triangle $C(T_D;E)\to C(A;E)\to Q\to C(T_D;E)[1]$, the $p$-component $e_p\in H^3_p(X;\mathrm{DR}_{\log})$ of the connecting morphism is $0$ when $(r_1,r_2)\ne(1,1)$.
\end{lem}
\begin{proof}[Sketch]
By Prop.~\ref{prop:moser} and Lemma~\ref{lem:res} we may assume $\sigma=c\,z_1z_2\partial_1\wedge\partial_2$ and $\nabla^{\log}e=(r_1\dlog z_1+r_2\dlog z_2)e$. All complexes are graded by the torus weight $\gamma\in\Z^2$, and the differentials preserve weights ($\lambda_i=\gamma_i+r_i$). Local cohomology can be computed with the Laurent tail $H^2_p(\mathcal O)=\mathrm{span}\{z^{-a}:a\ge1\}$. (a) $H^3_p(\mathrm{DR}_{\log})$ appears only in the weight with $\lambda_1=\lambda_2=0$ ($\gamma=-r$), and is non-zero only if $r_1,r_2$ are positive integers. (b) $H^2_p(Q)=\C$ sits in the single weight $\gamma=(-1,-1)$ (this requires an extra class $H^1_p(\mathcal O_D)\ni[c_1\ne c_2]$ in the local cohomology of $\mathcal O_D$), and $(-1,-1)=-r$ only if $r_1=r_2=1$. Since the connecting morphism preserves weight, $e_p=0$ for $(r_1,r_2)\ne(1,1)$. Together with the weight-by-weight Koszul computation there is no convergence problem (all maps are diagonal on monomials).
\tV\ Script \texttt{cp2/ncd\_local\_weights.py}: for weights $|\gamma_i|\le8$ and 78 pairs with $r_i\in\{-3..3,\pm\frac12,\frac13,-\frac23,\frac52\}$ (all $r_i\ne1$), the stalk cohomology $[0,0,1]$, the cone of local cohomology $=H^2_p$ (one class), and injectivity of $H(\mathrm{LOG})\to H(\mathrm{FULL})$ (connecting morphism $0$) were confirmed in exact arithmetic. 0 violations. For $r=(1,1)$ the connecting morphism is non-zero.
\end{proof}
\begin{proof}[Proof of Theorem~\ref{thm:main2}]
By Theorem~\ref{thm:main1} and Lemma~\ref{lem:stalk}, $C(T_D;E)\simeq\mathrm{DR}_{\log}\simeq\mathcal F$. By Lemmas~\ref{lem:Q} and \ref{lem:conn}, $C(A;E)\simeq\mathrm{DR}_{\log}\oplus\bigoplus_p\C_p[-2]$ in $D^b(\C_X)$. Apply $R\Gamma(X;\cdot)$.
(i), (ii): $\mathcal F=Rj_*L_U$ or $j_!L_U$, and since $X$ is compact $R\Gamma(X;j_!L_U)=R\Gamma_c(U;L_U)$.
(iii): $\chi(X;\mathcal F)=\sum_S\chi_c(S)\chi(\mathcal F_x)$, and the stalk Euler characteristic on $D$ is $0$ ($\wedge^*\C^s$ or $0$), so this equals $\chi_c(U)=\chi(U)=\chi(X)-\chi(D)$. Since $D$ is anticanonical, $\omega_D=\mathcal O_D$, and $\chi(\mathcal O_D)=0$ gives $\sum_j(2-2\tilde g_j)=2\mu$ ($\tilde D_j$ the normalizations), so $\chi_{\rm top}(D)=\mu$ and $\chi(U)+\mu=\chi(X)$. (Alternative proof in the simple normal crossings case: $\chi(U)=\int c_2(T_X(-\log D))$.)
\end{proof}
\begin{prop}[$r_i=1$: $E=K_X$]\tP\label{prop:K}
For $E=K_X$ we have $\wedge^bT\otimes K_X=\Omega^{2-b}$ and $E_1^{a,b}=H^{2-b,a}(X)$ (total degree $2+a-p$ with $p=2-b$). If $h^{p,q}(X)=0$ for $p\ne q$, then $E_1$ is concentrated in total degree $2$, the differentials (which raise total degree by 1) vanish, so $\HH^k=0$ ($k\ne2$) and $\dim\HH^2=\sum_ph^{p,p}=\chi(X)$ ($3$ for $\mathbb{CP}^2$, $4$ for $\mathbb{CP}^1\times\mathbb{CP}^1$).
\end{prop}

\subsection{Check on $\mathbb{CP}^1\times\mathbb{CP}^1$}
Let $\sigma=\Pi_F$, $F\in H^0(\mathcal O(2,2))$ ($-K=\mathcal O(2,2)$). \tP\ For smooth $D$, $\mathcal O(a,b)$ has a Poisson module structure only if $a=b=m$, the residue is $-m/2$, and the structure is unique (same argument as in \S\ref{sec:cp2}). \tV\ Exact computation for $m=-9,\dots,8$ (script \texttt{cp2/p1p1.py}):
\begin{center}\small
\begin{tabular}{lll}\toprule
$D$ & $m$ even $\ge0$ / even $\le-4$ & $m$ odd, $m=-2$\\\midrule
smooth $(2,2)$ curve & $[1,0,3,0,0]$ / $[0,0,3,0,1]$ & $[0,0,4,0,0]$\\
two transverse $(1,1)$ curves ($\mu=2$) & $[1,1,4,0,0]$ / $[0,0,4,1,1]$ & $[0,0,4,0,0]$\\
four lines $x_0x_1y_0y_1$ ($\mu=4$) & $[1,2,5,0,0]$ / $[0,0,5,2,1]$ & $[0,0,4,0,0]$\\
$(1,1)+(1,0)+(0,1)$ ($\mu=3$) & $[1,1,4,0,0]$ / $[0,0,4,1,1]$ & $[0,0,4,0,0]$\\\bottomrule
\end{tabular}
\end{center}
Period 2, Serre duality $m\leftrightarrow-m-4$, and $\chi=4$ always hold. For reducible $D$ the Poisson module structure is not unique ($H^0(\Omega^1(\log D))\ne0$), and the table is computed for the symmetric structure with equal residues on all components (Lemma~\ref{lem:res}).

\subsection{Verification of Theorem~\ref{thm:main2}}
\tV\ Theorem~\ref{thm:main2} and its closed formula (Theorem~\ref{thm:B} in \S\ref{sec:RR}) were compared with the exact computations above for a total of 196 values of $m$ in the following eight examples (script \texttt{cp2/ncd\_global\_check.py}): on $\mathbb{CP}^2$ a smooth cubic, a nodal cubic ($\mu=1$), a conic plus a secant line ($\mu=2$), a triangle ($\mu=3$) ($m=-15..15$), and the four $\mathbb{CP}^1\times\mathbb{CP}^1$ examples above ($m=-9..8$). 0 mismatches. The cuspidal cubic, conic plus tangent line (tacnode), three concurrent lines and non-reduced $D$ on $\mathbb{CP}^2$ are outside the scope of this theorem and are treated in \S\ref{sec:local}.


%--------------------------------------------------------------
\section{Serre-type theorems}\label{sec:serre}
%--------------------------------------------------------------
\subsection{Correspondence with the classical case}
The generalized counterparts of the ingredients of the classical Serre vanishing theorem (Theorem~\ref{thm:Serre}) are as follows (\tK\ the correspondence is the author's organization).
\begin{center}\small
\begin{tabular}{p{4.6cm}p{9.8cm}}\toprule
Classical & Counterpart in this note\\\midrule
Coherent sheaf $\mathcal F$ & Each term $K^b=\wedge^bT\otimes E$ of the Lie algebroid complex ($E$ a generalized holomorphic bundle). It is coherent.\\
Ample $\mathcal O(1)$ & An ample holomorphic line bundle $L_0$ which is also a Poisson module (generalized holomorphic). The Poisson structure is needed only so that $C(A;\cdot)$ is a complex; it plays no role in positivity.\\
Projective & $X$ projective (because $L_0$ exists).\\
$H^q(\mathcal F(m))=0\ (q>0)$ & $E_1^{a,b}=H^a(X,\wedge^bT\otimes E\otimes L_0^m)=0$ (\emph{$a>0$: the transversal (Dolbeault) degree}).\\
Conclusion: $H^*(\mathcal F(m))$ is concentrated in degree 0 & $\HH^*$ is the cohomology of the finite-dimensional complex of global sections $(\Gamma(\wedge^\bullet T\otimes E\otimes L_0^m),d)$.\\\bottomrule
\end{tabular}
\end{center}

\subsection{Serre-type theorem in any dimension}
\begin{thm}[Serre-type theorem ($\beta$-transforms)]\tP\label{thm:A1}
Let $X$ be a compact complex projective manifold ($\dim_\C X=n$), $\sigma$ any holomorphic Poisson structure ($[\sigma,\sigma]=0$; neither non-degeneracy nor a logarithmic condition is needed), $E$ a generalized holomorphic vector bundle (Poisson module), and $L_0$ an ample holomorphic line bundle with a Poisson module structure. Then there is $m_0$ depending only on $X,\sigma,E,L_0$ such that for $m\ge m_0$:
\begin{enumerate}
\item[(a)] $H^a(X,\wedge^bT\otimes E\otimes L_0^m)=0$ for all $a>0$ and all $b$;
\item[(b)] $\HH^k(E\otimes L_0^m)\cong H^k\bigl(\Gamma(X,\wedge^\bullet T\otimes E\otimes L_0^m),\,d\bigr)$. In particular $\HH^k=0$ for $k>n$;
\item[(c)] $\HH^0(E\otimes L_0^m)=$ the space of all flat (Casimir) holomorphic sections.
\end{enumerate}
\end{thm}
\begin{hyptab}
$X$ projective, $L_0$ ample & The hypotheses of classical Serre themselves. & (a): apply Hartshorne III.5.2 to the finitely many ($b=0,\dots,n$) coherent sheaves $\wedge^bT\otimes E$.\\\midrule
$L_0$ a Poisson module (and $E$) & $L_0^m$ gets the tensor-product module structure, so $C(A;E\otimes L_0^m)$ is a complex. & Definition of the complex $K^\bullet=C(A;E\otimes L_0^m)$. Unrelated to positivity.\\\midrule
$[\sigma,\sigma]=0$ & $d^2=0$. & Same.\\\midrule
(not needed) non-degeneracy of $\sigma$, normal crossings & --- & Not used in this theorem.
\end{hyptab}
\begin{proof}
$L_0^m$ has the tensor-product module structure (Leibniz rule). Applying Serre's theorem (Theorem~\ref{thm:Serre}, Hartshorne III.5.2) to the $n+1$ coherent sheaves $\wedge^bT\otimes E$ and taking $m_0$ the maximum of the thresholds gives (a). Each $K^b$ is $\Gamma$-acyclic, so for the bounded complex $K^\bullet$ (whose differential is $\C$-linear but not $\mathcal O$-linear) the spectral sequence $E_1^{a,b}=H^a(K^b)$ is concentrated at $a=0$ and $R\Gamma(K^\bullet)=\Gamma(K^\bullet)$ (acyclic resolution lemma). This is (b). Since $b\le n$, it vanishes for $k>n$. (c) is $H^0$ of the $\Gamma$-complex: $\HH^0=\ker(d_A:H^0(E)\to H^0(T\otimes E))$, i.e.\ flat sections.\qedhere
\end{proof}
\begin{rem}[Vanishing in total degree is false]\tP
For $\sigma=0$ (a complex structure), $\HH^k(L)=\bigoplus_{a+b=k}H^a(X,\wedge^bT\otimes L)$, so even for $m\gg0$ we have $H^0(\wedge^bT\otimes L^m)\ne0$, hence $\HH^b\ne0\ (b=0,\dots,n)$. For $m=0$ on $\mathbb{CP}^2$ this gives $[1,8,10,0,0]$ (\tV). Only the transversal degree $a>0$ and total degrees $>n$ vanish; ``Serre vanishing in total degree'' does not hold.
\end{rem}
\begin{rem}[\tV\ Effective thresholds]
By Bott's formula, the condition for $E_1^{a>0}=0$ is $m\ge-2$ on $\mathbb{CP}^2$ (for $m=-3$ the three terms $(2,0),(1,1),(0,2)$ are all non-zero), and $m\ge-1$ on $\mathbb{CP}^1\times\mathbb{CP}^1$ with $L_0=\mathcal O(1,1)$ ($m=-2$ is the exception $K_X$). Both agree with the range in which degrees $\ge3$ vanish in the computed tables.
\end{rem}

\subsection{Periodic stabilization}
\begin{thm}[Periodic stabilization (log-symplectic surfaces)]\tP\ (depends on Theorem~\ref{thm:main2} and its hypotheses)\label{thm:A2}
Let $X$ be a simply connected projective surface, $\sigma$ log-symplectic along $D$ (anticanonical, only ordinary double points), $L_0$ a line bundle with a Poisson module structure whose residues are $\rho_i=-a_i$ ($a_i\in\mathbb Q_{>0}$, constant along each component), and $U=X\setminus D$ affine ($D$ ample). Let $N$ be the least common multiple of the denominators of the $a_i$. If $m\ge1$ then
\[\HH^k(L_0^m)=0\ (k\ge3),\qquad\HH^*(L_0^m)=H^*(U;L_U^{(m)})\oplus\C^\mu[-2],\]
where $L_U^{(m)}$ is the local system with monodromy $e^{-2\pi ima_i}$ around $D_i$. In particular $\HH^*(L_0^m)$ depends only on $m\bmod N$.
\end{thm}
\begin{hyptab}
Residues $-ma_i<0$ & $r_i\ne1$ and not a positive integer. & Theorem~\ref{thm:main2}(i): $\mathcal F=Rj_*L_U$.\\\midrule
$U$ affine ($D$ ample) & Artin vanishing $H^k(U;L)=0\ (k>2)$. & Vanishing in degrees $\ge3$. (Theorem~\ref{thm:A1}(b) gives the same conclusion from projectivity and ampleness and makes affineness unnecessary.)\\\midrule
$X$ simply connected & The abelianization of $\pi_1(U)$ is generated by meridians. & The local system is determined by $m\bmod N$.\\\midrule
$a_i\in\mathbb Q_{>0}$ & $c_1(L_0)$ is a positive rational combination of the boundary components (``positive with respect to the boundary''). & Periodicity.
\end{hyptab}
\tV\ Period 3 for $\mathbb{CP}^2$ ($L_0=\mathcal O(1)$, $a=1/3$), period 2 for $\mathbb{CP}^1\times\mathbb{CP}^1$ ($L_0=\mathcal O(1,1)$, $a=1/2$), for $m\ge1$. Confirmed for all eight computed examples.

\begin{conj}[General case]\tK\label{conj:A3}
For a general generalized complex manifold $M$ and a generalized holomorphic line bundle $L_0$ that is ``positive'' in a suitable sense (candidate: ample on the leaf space of the regular locus, and with positive residues along the locus where the type jumps), $\HH^*(L_0^m)$ vanishes in total degrees $>n_{\max}$ for $m\gg0$ and is periodic in $m$. For a general $\Jc$ that is neither a $\beta$-transform nor of regular type, there is no way to define the ``$E_1$ term of a coherent filtration $K^b$'', and this is the main obstruction to going beyond Theorem~\ref{thm:A1}.
\end{conj}

\subsection{Relation to the theorem of Pal--Poddar}
\begin{prop}[Consistency in the product case]\tP\label{prop:PPprod}
Let $M=Y\times X$ ($Y$ a compact symplectic manifold, $X$ a projective complex manifold) and $E=\mathrm{pr}_Y^*F\boxtimes\mathrm{pr}_X^*G_0^m$ ($F$ a flat bundle on $Y$, $G_0$ ample). Then the generalized Dolbeault complex is the tensor product of the de Rham complex of $Y$ and a complex on $X$, and
\[\HH^k=\bigoplus_{j+b+a=k}H^j_{dR}(Y;F)\otimes H^a(X,\wedge^bT_X\otimes G_0^m)\]
(K\"unneth). By classical Serre the terms with $a>0$ vanish for $m\gg0$, and the remaining $H^j_{dR}(Y;F)\otimes H^0(X,\wedge^bT_X\otimes G_0^m)$ are in general non-zero up to total degree $\dim Y+\dim_\C X$. Hence Serre-type vanishing for $\mathrm{pr}_X^*G_0^m$ holds only for the transversal degree $a$.
\end{prop}
\src{Pal--Poddar, \emph{Strong generalized holomorphic principal bundles}, arXiv:2404.18113, \S11, Thm.~11.12: if $M$ is a compact generalized complex manifold of regular type $k$ with $M/\mathcal S$ a K\"ahler orbifold and $E$ is a positive SGH line bundle, then (1) $H^q(M,(\mathcal G^*M)^p\otimes E)=0$ for $p+q>k$, and (2) for every SGH vector bundle $E'$ there is $m_0$ such that $H^q(M,E'\otimes E^m)=0$ for $m\ge m_0$, $q>0$.}
\begin{rem}[\tQ]
The groups $H^q(M,(\mathcal G^*M)^p\otimes E)$ of Pal--Poddar (which by Cor.~8.10 are $H^{p,q}_{d_L}(M,E)$) are sheaf cohomology groups of transversal forms that are constant along the leaves, and $q$ is the transversal (complex) degree (confirmed by re-reading \S\S5--11, 13; the extraction went through summaries, so details of formulas are unverified). Hence the reading of the proposition above (vanishing only in the transversal degree $a$) is consistent with Thm.~11.12. The groups are not the same, however: Pal--Poddar use exterior powers of $\mathcal G^*M$, this note tensors with $\wedge^bT$ (tangent side), and moreover the $\HH$ of this note is the total complex including the leafwise de Rham factor. The proposition above is a computation in the product case and is not itself a verification of Thm.~11.12. Theorem~\ref{thm:A1} does not assume regular type or an orbifold leaf space, and includes type-jumping $\beta$-transforms (e.g.\ types $0$ and $n$ coexisting), at the price of requiring the underlying complex manifold to be projective (in the $\beta$-transform case only).
\end{rem}

%--------------------------------------------------------------
\section{Duality}\label{sec:duality}
%--------------------------------------------------------------
\begin{thm}[Serre duality for generalized holomorphic line bundles]\tP\label{thm:dual}
Let $X$ be any compact complex manifold ($\dim_\C X=n$), $\sigma$ any holomorphic Poisson structure (no K\"ahler, projective, log-symplectic or non-degeneracy assumption is needed; $\sigma=0$ is included), and $L$ a line bundle that is a Poisson module. Then for all $k$
\[\HH^k(L)^*\;\cong\;\HH^{2n-k}\bigl(L^\vee\otimes K_X^{2}\bigr).\]
$L^\vee\otimes K_X^2$ carries the tensor-product module structure (frame $e'$ with $ee'=\omega^2$ and $V'=-V+2V_K=-V-2M$; the canonical module structure on $K_X$ is $V_K^k=\sum_i\partial_i\sigma^{ki}$, and $M^j=\sum_i\partial_i\sigma^{ij}$ is the modular vector field with respect to $\omega=dz_1\wedge\cdots\wedge dz_n$).
\end{thm}
\begin{hyptab}
$X$ compact & Closed range and finite-dimensionality of elliptic complexes; no boundary terms in Stokes. & Step (4) (identification of the dual).\\\midrule
$[\sigma,\sigma]=0$, $L$ a Poisson module & $D^2=0$, $D'^2=0$ ($L'$ is also a module). & To treat them as complexes.\\\midrule
(not needed) K\"ahler / non-degenerate / logarithmic & Ellipticity is checked directly from the principal symbol, and the transposition is an identity of local operators, so no global primitive is needed. & Steps (1), (3).
\end{hyptab}
\begin{proof}[Proof (sketch)]
(0) Dolbeault model: $\mathrm{Tot}^k(L)=\bigoplus_{a+b=k}A^{0,a}(\wedge^bT\otimes L)$, $D=\dbar+d_K$ ($d_K(\phi\otimes Pe)=(-1)^a\phi\otimes d(Pe)$, $d=[\sigma,\cdot]+V\wedge$). $D^2=0$ follows from $\dbar^2=0$, $d^2=0$, $\dbar d_K+d_K\dbar=0$ (Prop.~\ref{prop:ident}(b)).
(1) Ellipticity (\tP): at a real cotangent vector $\xi\ne0$, the principal symbol of $D$ is the exterior product with $v(\xi)=(\xi^{0,1},\sigma^\sh\xi^{1,0})\in L^*=T^{*0,1}\oplus T^{1,0}$, and $(\wedge L^*,v\wedge)$ is exact if and only if $v\ne0$. If $v=0$ then $\xi^{0,1}=0$, and since $\xi$ is real, $\xi=0$. Hence the complex is elliptic at every point, including points where $\sigma$ degenerates (this is a direct verification, in the present case, of Gualtieri's Prop.~\ref{prop:ellip}). The same holds for $L'$.
(2) Perfect pairing: integrating $E^b(L)\times E^{n-b}(L')\to\wedge^nT\otimes K^2=K$ and $A^{0,a}\times A^{0,n-a}\to A^{0,n}(K)=A^{n,n}$ gives $\mathrm{Tot}^k(L)\times\mathrm{Tot}^{2n-k}(L')\to\C$ (a perfect pairing at each point).
(3) Transposition: from the following two identities, the formal transpose $D^t$ equals $-(-1)^{\deg}D'$.
\end{proof}
\begin{lem}[Lemma D: top-degree part of the Schouten bracket]\tP\ \tV\label{lem:D}
For $T=\sum_j(-1)^ju_j\,\partial_1\wedge\cdots\widehat{\partial_j}\cdots\wedge\partial_n\in\wedge^{n-1}T$ we have $[\sigma,T]_{\rm top}=-\sum_{i,j}\partial_i(\sigma^{ij}u_j)+2\sum_jM^ju_j$ and $(M\wedge T)_{\rm top}=\sum_jM^ju_j$.
\end{lem}
\begin{lem}[Lemma E: transposition identity]\tP\ \tV\label{lem:E}
For $P\in\wedge^bT$, $Q\in\wedge^{n-b-1}T$ (smooth coefficients), $d=[\sigma,\cdot]+V\wedge$, $d'=[\sigma,\cdot]+V'\wedge$, $V'=-V-2M$, and $\langle P,Q\rangle:=(P\wedge Q)_{\rm top}$, we have
\[\langle dP,Q\rangle+(-1)^b\langle P,d'Q\rangle=-\sum_i\partial_i\Bigl(\sum_j\sigma^{ij}u_j\Bigr),\quad T=P\wedge Q=\sum_j(-1)^ju_j\partial_{\hat j}.\]
Proof: the Leibniz rule $[\sigma,P\wedge Q]=[\sigma,P]\wedge Q+(-1)^bP\wedge[\sigma,Q]$ ($\sigma$ has even degree); the $V$-terms give $V\wedge P\wedge Q+(-1)^bP\wedge(-V)\wedge Q=0$, and the $M$-term is $-2M\wedge T$. The result follows from Lemma~D.
\tV\ For $n=2,3,4$, with random polynomial and Laurent monomial data and arbitrary bivector fields $\sigma$ (the identity itself holds even if $\sigma$ is not Poisson), the identity was confirmed in exact arithmetic (\texttt{cp2/duality\_ndim\_exact.py}). The constant $c=-2$ and the sign $(-1)^b$ were found by solving a system of equations independently of the derivation (\texttt{cp2/duality\_ndim.py}). $[\sigma,M]=0$ and $d^2=0$ ($V=M$, $\sigma$ Poisson) were also confirmed.
\end{lem}
\begin{proof}[Continuation of the proof of Theorem~\ref{thm:dual}]
(3) continued: the right-hand side of Lemma~E is an exact form of the type $\partial_iF^i\,\omega\wedge d\bar z^{1\cdots n}=\pm d(F^i\iota_{\partial_i}\omega\wedge d\bar z^{1\cdots n})$ (the $\dbar$ component vanishes against $d\bar z^{1\cdots n}$), so by Stokes its integral is $0$. The $\dbar$ part is the classical Serre duality computation: $\langle\dbar\alpha,\beta\rangle=-(-1)^{a+b}\langle\alpha,\dbar\beta\rangle$. The $d$ part is $\langle d_K\alpha,\beta\rangle=-(-1)^{a+b}\langle\alpha,d'_K\beta\rangle$. That the two signs agree was checked by hand (not machine-checked \tQ). Since $D^t$ is a local operator, the equality in coordinate charts gives globally $\langle D\alpha,\beta\rangle=-(-1)^{\deg\alpha}\langle\alpha,D'\beta\rangle$.
(4) Conclusion: $D,D'$ are elliptic complexes on a closed manifold, transposes of each other (up to sign) with respect to the perfect pairing. By Hodge theory for elliptic complexes, $\HH^k$ is finite-dimensional with closed range, $\HH^k(L)^*$ is the cohomology of the transposed complex (with distributional coefficients), and by elliptic regularity it equals the one with smooth coefficients, $\HH^{2n-k}(L')$.\qedhere
\end{proof}
\begin{rem}[Consistency checks\ \tV]
(i) $\sigma=0$: classical Serre duality (using $\wedge^bT^*=\wedge^{n-b}T\otimes K$). (ii) $\mathbb{CP}^2$, $K^2=\mathcal O(-6)$: $m\leftrightarrow-m-6$; $\mathbb{CP}^1\times\mathbb{CP}^1$, $K^2=\mathcal O(-4,-4)$: $m\leftrightarrow-m-4$ (consistent with the tables in \S\ref{sec:cp2} and \S\ref{sec:main2}). (iii) In the log-symplectic case the residue of $K_X$ is $1$ ($r=1$ in Prop.~\ref{prop:K}), so the residue of $L'$ is $2-r_i$; this is the origin of the symmetry $r\leftrightarrow2-r$ in the formula of the main theorem (in any even dimension).
\end{rem}


%--------------------------------------------------------------
\section{Riemann--Roch-type formulas}\label{sec:RR}
%--------------------------------------------------------------
\subsection{The total Euler characteristic is a topological invariant}
\begin{prop}\tP\label{prop:euler}
Let $X$ be a compact complex manifold, $\sigma$ any holomorphic Poisson structure, and $E$ a Poisson module of rank $\rk E$. Then
\[\sum_k(-1)^k\dim\HH^k(E)=\sum_b(-1)^b\chi(X,\wedge^bT\otimes E)=(-1)^n\,\rk E\cdot\chi(X).\]
\end{prop}
\begin{proof}
The complex $C(A;E)=(\wedge^bT\otimes E,d)$ has a filtration by $b$ (the stupid filtration), and the spectral sequence $E_1^{a,b}=H^a(X,\wedge^bT\otimes E)$ is finite-dimensional and converges to $\HH^{a+b}$. The Euler characteristic equals that of the $E_1$ term, and $\sum_b(-1)^b\chi(\wedge^bT\otimes E)$ is, by Hirzebruch--Riemann--Roch (Theorem~\ref{thm:HRR}),
$\int\ch(E)\Td(T)\sum_b(-1)^b\ch(\wedge^bT)$. Since $\Td(T)=\prod\frac{x_i}{1-e^{-x_i}}$ and $\sum_b(-1)^b\ch(\wedge^bT)=\prod(1-e^{x_i})$, the product is $\prod(-x_ie^{x_i})=(-1)^ne^{c_1}c_n(T)$. As $c_n$ is of top degree, only the rank of $\ch(E)$ contributes, giving $(-1)^n\rk E\int c_n(T)=(-1)^n\rk E\,\chi(X)$.\qedhere
\end{proof}
\begin{rem}[\tP\ Meaning]
For $\mathbb{P}^1$ and $E=\mathcal O(m)$ the formula reads $\chi(\mathcal O(m))-\chi(\mathcal O(m+2))=(m+1)-(m+3)=-2$ (independent of $m$). That is, the Riemann--Roch formula for the \emph{total} Euler characteristic of the generalized Dolbeault complex is a topological formula that does not see $c_1(E)$, and is not a generalization of the classical $\chi(M,E)=\int\ch(E)\Td(M)$. The classical $\chi(M,E)$ corresponds to the Euler characteristic of the ``column'' $b=0$, $H^\bullet(X,E)$. For a general $\Jc$ this column is not preserved by $d$, and some refinement is needed.
\end{rem}

\subsection{Closed formula for the Poincar\'e polynomial}
\begin{thm}[Riemann--Roch-type closed formula (surfaces)]\tP\ (depends on Theorem~\ref{thm:main2} and the hypotheses below)\label{thm:B}
In the setting of Theorem~\ref{thm:main2} assume in addition: (H1) $H^1(X;\C)=0$ (for $b_1(U)=c-\rho$); (H2) $U$ is affine (for vanishing in degrees $\ge3$); (H3) all components have the same residue $r_i=r\ne1$; (H4) if $r\notin\Z$, then $H^0(U;L_U)=H^1(U;L_U)=0$ (G2: see the remark below). Put $b_1=c-\rho$ ($c$ the number of components, $\rho$ the dimension of the span of $[D_i]\in H^2(X;\C)$), $\chi_U=\chi(X)-\mu$, $b_2=\chi_U-1+b_1$. Then
\begin{align*}
P_E(t)&=1+b_1t+(b_2+\mu)t^2 &&(r\in\Z_{\le0}),\\
P_E(t)&=(b_2+\mu)t^2+b_1t^3+t^4 &&(r\in\Z_{\ge2}),\\
P_E(t)&=\chi(X)\,t^2 &&(r\notin\Z).
\end{align*}
Equivalently: $P_E(t)=\chi(X)t^2+[r\le0]\,\Delta(t)+[r\ge2]\,t^4\Delta(1/t)$, $\Delta(t)=(1+t)(1+(b_1-1)t)$.
\end{thm}
\begin{hyptab}
(H1) $H^1(X;\C)=0$ & From the Gysin sequence $H^0(D)\to H^2(X)$, $H^1(U)=\ker(\C^c\to H^2(X))$. & $b_1(U)=c-\rho$.\\\midrule
(H2) $U$ affine & Artin vanishing. & $b_3=b_4=0$, hence $b_2=\chi_U-1+b_1$.\\\midrule
(H3) equal residues & $\mathcal F$ is either $Rj_*$ or $j_!$. & Theorem~\ref{thm:main2}(i), (ii).\\\midrule
(H4) $H^0(U;L)=H^1(U;L)=0$ for $r\notin\Z$ & So that $H^*(U;L)=[0,0,\chi_U]$. & The row $r\notin\Z$. For nodal curves on $\mathbb{CP}^2$, $\pi_1(U)$ is abelian (Fulton, Ann.\ Math.\ 111 (1980) 407--409; Deligne, S\'em.\ Bourbaki 543). \tC\\\midrule
$r_i\ne1$ & Hypothesis of Theorem~\ref{thm:main2}. & Lemma~\ref{lem:Q}.
\end{hyptab}
\begin{proof}
Substitute into Theorem~\ref{thm:main2}(i), (ii): $H^*(U)=[1,b_1,b_2,0,0]$ ($r\in\Z_{\le0}$, trivial local system), $H_c^*(U)$ (by Poincar\'e duality $[0,0,b_2,b_1,1]$, $r\in\Z_{\ge2}$), $H^*(U;L)=[0,0,\chi_U,0,0]$ ($r\notin\Z$: from (H2), (H4) and $\chi(U;L)=\chi(U)$), and add $\mu t^2$. The equivalent form uses $\chi(X)=\chi_U+\mu=1-b_1+b_2+\mu$.
\end{proof}
\begin{rem}[Properties\ \tP]
(i) $P_E(-1)=\chi(X)$ for all $r$ (consistent with Prop.~\ref{prop:euler}, $n=2$). (ii) The duality $P_{2-r}(t)=t^4P_r(1/t)$ is evident from the formula (a formula-level consequence of Theorem~\ref{thm:dual}). (iii) $\mu$ and $b_2$ are absorbed into $\chi(X)t^2$, and the jump term $\Delta$ depends only on $b_1$ (the number of linear relations among the $[D_i]$). (iv) For $E=L_0^m$, the quasi-polynomial in $m$ that classical Riemann--Roch would give is the constant $\chi(X)$, and the dependence on $m$ appears only through the integrality (period $N$) and the sign of the residue $r=-ma$.
\end{rem}
\begin{rem}[\tQ\ On G2]
(H4) was checked in the eight verified examples (four curves on $\mathbb{CP}^2$ and four on $\mathbb{CP}^1\times\mathbb{CP}^1$). For the smooth $(2,2)$ curve and the two $(1,1)$ curves on $\mathbb{CP}^1\times\mathbb{CP}^1$ we cited Timmerscheidt (\emph{Mixed Hodge theory for unitary local systems}, J.\ reine angew.\ Math.\ 379 (1987) 152--171) and Arapura (\emph{Geometry of cohomology support loci for local systems I}, arXiv:alg-geom/9612006, J.\ Algebraic Geom.\ 6 (1997) 563--597 (volume/pages from memory \tQ), the torsion-translate theorem) (applicability not fully checked \tQ). In general $H^1(U;L_\chi)=0\Leftrightarrow\chi\notin V^1(U)$ (the first characteristic variety), which can fail for non-nodal $D$ (e.g.\ three concurrent lines) (Libgober, Arapura).
\end{rem}
\tV\ The closed formula agrees with the exact computations for a total of 196 values: $\mathbb{CP}^2$ (smooth, nodal, conic plus secant, triangle; $m=-15..15$) and $\mathbb{CP}^1\times\mathbb{CP}^1$ (4 examples, $m=-9..8$) (0 mismatches, \texttt{cp2/ncd\_global\_check.py}). The formula-level symmetry $r\leftrightarrow2-r$ was checked on a grid of 73 rational values of $r$ (\texttt{cp2/serre\_rr\_check.py}).

\subsection{Index-theoretic outlook}\label{sec:index}
The content of this subsection is conjectural; no index theorem is proved.
\begin{obs}[\tK\ Indicial-root picture (qualitative)]
At smooth points the logarithmic complex acts on the monomial $z^ne$ by the factor $(n-r)$ (Lemma~\ref{lem:stalk}). An indicial root exists exactly when $r\in\Z$; for $r\notin\Z$ the weight-graded complexes are acyclic and $P_r=\chi(X)t^2$. For $r\in\Z$, what remains is a subquotient of the cohomology of the boundary $\partial U$ of a tubular neighbourhood of $D$ (via the Gysin sequence, $H^*(D)\otimes(1\oplus t)$), and since $\chi(\partial U)=0$ it is invisible in the Euler characteristic. This is qualitatively consistent with Melrose's APS / b-calculus picture (Melrose, \emph{The Atiyah--Patodi--Singer Index Theorem}, A K Peters 1993).
\end{obs}
\begin{conj}[Index-type Riemann--Roch (log-symplectic surfaces)]\tK\label{conj:I3}
(a) $\sum(-1)^k\dim\HH^k(L)=\int_Xc_2(T_X(-\log D))+\mu=\chi(X)$ is given by an Atiyah--Singer-type integral without boundary term. (b) The jump term $J(r)=P_r(t)-\chi(X)t^2$ is non-zero only for $r\in\Z\setminus\{1\}$, equal to $\Delta(t)$ for $r\le0$ and to $t^4\Delta(1/t)$ for $r\ge2$. $\Delta(t)$ is the Poincar\'e polynomial of a certain subquotient of the kernel $H^*(\partial U)$ of the indicial roots (fixed by the weight filtration of the long exact sequence $H_c^*(U)\to H^*(U)\to H^*(\partial U)$). (c) In general dimension and for non-normal-crossing $D$, the jump term contains local terms governed by the spectral numbers of the singularities (\S\ref{sec:local}).
\end{conj}
\begin{rem}[\tC\ State of the literature and obstacles]
\begin{itemize}
\item \emph{Analogy with the $b$-index}: Braverman--Loizides--Song, \emph{Geometric quantization of $b$-symplectic manifolds}, arXiv:1910.10016 shows that the quantization index of a $b$-symplectic manifold is given by $\int\mathrm{Td}({}^bTM)\mathrm{Ch}(L)$ and that the boundary term vanishes in the limit (title and authors checked \tC). This is analogous to the fact that the Euler characteristic here is of the form $\chi(X)=\int c_2(T_X(-\log D))+\mu$, independent of $r$.
\item \emph{Index on Lie groupoids}: Pflaum--Posthuma--Tang, \emph{The index of geometric operators on Lie groupoids}, arXiv:1308.0236 (Indag.\ Math.\ 25 (2014) 1--18); Qiao--So, arXiv:2304.13345 (boundary groupoids; no eta term in codimension $\ge2$); Ammann--Lauter--Nistor, arXiv:math/0304044; Mazzeo--Melrose, arXiv:math/9812120; Hausel--Hunsicker--Mazzeo, arXiv:math/0207169 (Duke Math.\ J.\ 122 (2004), pages \tQ). Titles and authors were checked; applicability to the present problem was not.
\item \emph{Poisson cohomology}: Lanius, arXiv:1605.03854 ($C^\infty$ Poisson cohomology of log-symplectic manifolds). This is a different object from the holomorphic Poisson hypercohomology of this note; there is only an analogy in the structure ``$b$-part $+$ boundary term''.
\item \emph{What was not found}: an index theorem for the complex Lie algebroid $L=T^{0,1}\bowtie A^{1,0}$, or a $b$-Dolbeault index with the residue as a weight. The real Lie algebroid $T_X(-\log D)$ is integrable for smooth $D$ by Gualtieri--Li, arXiv:1206.3674. \tC\ The normal crossing case was not checked.
\item \emph{Unimodularity}: the Pflaum--Posthuma--Tang trace requires an invariant density, which for Poisson structures corresponds to the vanishing of the modular class. In the log-symplectic case the modular vector field is expected to vanish thanks to the logarithmic volume form, but this was not checked \tQ.
\end{itemize}
\end{rem}

%--------------------------------------------------------------
\section{The toric case and higher dimensions}\label{sec:toric}
%--------------------------------------------------------------
\subsection{Toric log-symplectic structures}
Let $X$ be a smooth projective toric variety ($N=\Z^n$, $n=2k$) with rays $u_i$ ($i\in I$) and $D$ the toric boundary (simple normal crossings). With respect to the logarithmic frame $E_j=z_j\partial_j$, a constant non-degenerate skew-symmetric matrix $A$ gives $\sigma=\tfrac12\sum A_{ij}E_i\wedge E_j$ (torus-invariant, log-symplectic along $D$). Give $L=\mathcal O(\sum a_iD_i)$ the flat logarithmic connection $d+\rho$ ($\rho\in M_\C$). The residues are $r_i=\langle\rho,u_i\rangle-a_i$, and $c_1(L)=-\sum r_i[D_i]$. \tP\ Every module is of this form ($H^0(\Omega^1(\log D))=M_\C$, Lemma~\ref{lem:res}).

\begin{thm}[Weight formula]\tP\ \tV\label{thm:weight}
Assume $E_1^{a,b}=H^a(X,\wedge^bT\otimes L)=0\ (a>0)$ (e.g.\ in the range of Theorem~\ref{thm:A1}). For $\chi\in M$ put $c_i(\chi)=\langle\chi,u_i\rangle+a_i$, $C(\chi)=\{i:c_i(\chi)=-1\}$, $W(\chi)=\mathrm{span}(u_i:i\in C(\chi))$, $s=\dim W$. Then
\[P_{\HH}(t)=\sum_{\chi\in\Lambda}t^{s(\chi)}(1+t)^{n-s(\chi)},\quad\Lambda=\{\chi\in M:c_i(\chi)\ge-1\ \forall i,\ A(\chi+\rho)\in W(\chi)\}.\]
\end{thm}
\begin{hyptab}
Toric, torus-invariant $\sigma$, constant $A$ & The complex of global sections decomposes by torus characters $\chi\in M$ ($T$-weights), and each weight is a finite-dimensional linear algebra problem. & Proof: determination of the weight-$\chi$ subcomplex $G^b_\chi$.\\\midrule
$E_1^{a>0}=0$ & $\HH^*$ equals the cohomology of the complex of global sections. & Theorem~\ref{thm:A1}(b).\\\midrule
Non-degenerate $A$ & $\sigma^\sh:B\cong T_D$ holds, and $\wedge^b\sigma^\sh$ sends the logarithmic de Rham differential $(\chi+\rho)\wedge$ to $v_\chi\wedge$ ($v_\chi=A(\chi+\rho)$). & Lemma~\ref{lem:iso}.
\end{hyptab}
\begin{proof}[Sketch of proof]
The condition for $z^\chi w$ ($w\in\wedge^bN_\C$) to be a holomorphic section of $\wedge^bT\otimes L$ is: $c_i(\chi)\ge-1$ along each $D_i$, and if $c_i(\chi)=-1$ then $u_i\wedge w=0$. So the weight-$\chi$ part is $G^b_\chi=\{w\in\wedge^bN_\C:u_i\wedge w=0\ (i\in C(\chi))\}$ (empty if some $c_i<-1$), and $G^\bullet_\chi\cong\wedge^{\bullet-s}(N/W)\otimes\wedge^sW$. The differential is $v_\chi\wedge$. If $v_\chi\bmod W=0$ the differential is $0$ and the Poincar\'e polynomial is $t^s(1+t)^{n-s}$; otherwise the complex is acyclic.
\end{proof}
\tV\ The closed formula agrees with a direct linear-algebra computation (\texttt{cp2/toric\_hh.py}) in 60 random cases ($\mathbb P^2$, $\mathbb P^1\times\mathbb P^1$, $\mathbb P^4$, $\mathbb P^2\times\mathbb P^2$, $(\mathbb P^1)^4$).

\begin{cor}\tP
(a) If $E_1^{a>0}=0$, the total Euler characteristic is the number of weights with $s(\chi)=n$ (the number of vertices) $=\chi(X)$. (Consistent with Prop.~\ref{prop:euler}.) (b) Origin of the contributions: $s=0$ means $A(\chi+\rho)=0$, i.e.\ $\chi=-\rho$, and if $\rho\in M$ and all $r_i\le0$ it contributes $(1+t)^n=H^*(U)$ (the $Rj_*$ case). $s=n$ always satisfies the condition and gives the vertex weights, one $t^n$ each. For $0<s<n$ one needs $x=\chi+\rho\in A^{-1}W$ and $\langle x,u_i\rangle=r_i-1\ (i\in C)$, which has solutions only at ``resonances'' on the integral lattice. (c) For $s=n-1$ and all $r_i=1\ (i\in C)$, the condition holds automatically and classes $t^{n-1}(1+t)$ appear, in number equal to the number of lattice points on the toric curve $C$, namely $\deg(L|_C)+1$ (corresponding to the resonance $r_i=1$ in dimension 2).
\end{cor}

\subsection{Generic $\sigma$ and dimension 4}
\begin{obs}[\tV]
For generic $A$ (entries that are integers of very large size), $\Lambda=\{\text{vertices}\}\cup\{\text{resonances}\}\cup\{-\rho\}$ ($-\rho$ is included when $\rho\in M$ and all $r_i\le1$), and $P_{\HH}=$ ($U$-part) $+\#\text{vertices}\cdot t^n+$ (resonance classes for $r=1$). The $U$-part is $(1+t)^n$ if $\rho\in M$ and all $r_i\le0$, and $0$ otherwise. In exact arithmetic modulo the prime $2^{61}-1$, there were 0 mismatches in 3626 cases ($\mathbb P^4$ ($m=1,2$), $\mathbb P^2\times\mathbb P^2$, $(\mathbb P^1)^4$, $\mathbb P^1\times\mathbb P^3$, $\mathbb P^1\times\mathbb P^1\times\mathbb P^2$, $\rho\in\{-2,..,1\}^4$ and three non-integral $\rho$). Examples: $\mathbb P^4$, $m=1$: $[1,4,6,4,6]$; $\mathbb P^2\times\mathbb P^2$: $[1,4,6,4,10]$.
\end{obs}
\begin{obs}[New phenomenon: dependence on the arithmetic of $\sigma$\ \tV]\label{obs:arith}
In dimension 4, $\HH^*$ is not determined by $(X,D,L,\text{residues})$ alone. For the product structure $\sigma$ on $\mathbb P^2\times\mathbb P^2$ (block-diagonal $A$) it is $[1,4,12,16,16]$ (the product of the results for two surfaces, K\"unneth), and for generic $A$ it is $[1,4,6,4,10]$. The difference $(0,0,6,12,6)=6t^2(1+t)^2$ is $t^2H^*(O_\tau)$ for six codimension-2 strata $\{\text{double point of one factor}\}\times(\mathbb C^*)^2$. Mechanism: on a stratum with $|C|=2$, if $u_i,u_j$ are symplectic with respect to $\sigma$, then $x_C(A,r)\in A^{-1}W$ is uniquely determined, and it contributes only under the arithmetic condition that this is a lattice translate of $\rho$. \tK\ $\HH$ jumps up on countably many proper analytic subsets of the parameter space of log-symplectic structures (upper semicontinuity of generalized Dolbeault cohomology), and the generic value is the minimum.
\end{obs}
\begin{conj}[General simple normal crossings $D$]\tK
For a general (non-toric) simple-normal-crossings log-symplectic $(X^{2k},D)$, general $\sigma$, and $L$ with constant residues $r_i\ne1$: $\HH^*(L)=H^*(X;\mathcal F)\oplus\bigoplus_{p\in D_{(n)}}\C[-n]$ ($D_{(n)}$ the set of points where $n$ components meet). Consistency with the Euler characteristic holds only when the $\chi$ of the lower-dimensional strata $D_k^\circ$ vanish (true in the toric case, unproved in general). The computations of \S\ref{sec:toric} are for the toric case only.
\end{conj}

%--------------------------------------------------------------
\section{Local theory of non-nodal singularities}\label{sec:local}
%--------------------------------------------------------------
When $D$ has cusps, tacnodes, triple points etc., Theorem~\ref{thm:main2} does not apply. What follows are computational observations and conjectures. Local model: $\sigma=h\,\partial_z\wedge\partial_w$, module $\mathcal O\,h^s$, $d_0f=(hf_w+sfh_w)\partial_z-(hf_z+sfh_z)\partial_w$, $d_1(p,q)=h(p_z+q_w)-(1-s)(h_zp+h_wq)$. Globally $s=r=-m/3$ ($\mathbb{CP}^2$), $-m/2$ ($\mathbb{CP}^1\times\mathbb{CP}^1$).

\subsection{Stalk cohomology}
\tV\ For germs of quasi-homogeneous isolated singularities, the stalk dimensions $(h^0,h^1,h^2)$ were computed by weight-graded exact linear algebra (script \texttt{local\_stalk.py} etc.) (assuming that the cohomology of the formal germ equals that of the convergent germ):
\begin{center}\small
\begin{tabular}{lccc}\toprule
Singularity & generic $s$ & $s$ an integer $\le0$ & $s$ a spectral number mod $\Z$, $s<\mathrm{lct}$\\\midrule
smooth ($z$) & $(0,0,0)$ & $(1,1,0)$ & ---\\
double point ($zw$) & $(0,0,1)$ & $(1,2,2)$ & ---\\
cusp ($z^2-w^3$) & $(0,0,2)$ & $(1,1,2)$ & $(0,1,3)$\\
tacnode ($z^2-w^4$) & $(0,0,3)$ & $(1,2,4)$ & $(0,1,4)$\\
$A_4$ ($z^2-w^5$) & $(0,0,4)$ & $(1,1,4)$ & ---\\
$D_4$ ($zw(z-w)$) & $(0,0,4)$ & $(1,3,6)$ & $(0,1,5)$\\
$D_5$ ($w(z^2-w^3)$) & $(0,0,5)$ & $(1,2,6)$ & $(0,1,6)$\\\bottomrule
\end{tabular}
\end{center}
\begin{obs}[\tV\ Empirical rule]
In terms of the spectral numbers (lct is the smallest spectral number: smooth and double point $1$, cusp $5/6$, tacnode $3/4$, $D_4$ $2/3$, $D_5$ $5/8$), for $s\ne1$,
\[\text{stalk}_p=\C^{\mu_p}[-2]\ \oplus\ \begin{cases}\text{stalk}(Rj_*L_s)_p&(s<\mathrm{lct}_p)\\0&(s>\mathrm{lct}_p)\end{cases}\]
($\mu_p$ is the Milnor number $=$ Tjurina number). The stalk of $Rj_*$ is governed by the monodromy of the Milnor fibre: one class when $s\equiv\alpha_j\bmod\Z$ ($\alpha_j$ a spectral number), and a contribution to $H^0,H^1$ when $s\in\Z$. All 7 types, $s\in\frac1{12}\Z$ or $\frac18\Z$, $|s|\le5$: no contradiction. The explanation via Milnor monodromy is a consistency argument, not a proof. For $s=1$ (the twist by $K_X$) the stalk is infinite-dimensional ($d_Q=0$), but globally it can be handled by the $E_1$-term argument (Prop.~\ref{prop:K}).
\end{obs}
\subsection{The triple-point anomaly and a global conjecture}
\tV\ On $\mathbb{CP}^2$ with $xy(x-y)$ (a triple point): $m=1,2,4,5,\ldots$ and $m=-1$ give $[0,1,4,0,0]$, $m=-5,-7,-8$ give $[0,0,4,1,0]$, and $m=-2,-3,-4$ give $[0,0,3,0,0]$. By the empirical rule above, with the $D_4$ spectral numbers $\{2/3,1,1,4/3\}$, $\mathrm{lct}=2/3$, $s=-m/3$: for $m\ge-1$, $s<2/3$ adds $H^*(U;L)=(0,1,0)$ giving $[0,1,4,0,0]$; for $m\le-5$, $s>2/3$ adds $H_c^*=(0,0,0,1,0)$ giving $[0,0,4,1,0]$; the anomalous $m=-2,-4$ have $s=2/3,4/3$, where one class of $\C^4$ (degree 2) and one class of $H_c$ (degree 3) cancel (the $d_3$ of the spectral sequence $H^p(X,\mathcal H^q)\Rightarrow\HH$). \tK\ This explanation is conjectural. The spectral numbers of the cusp ($5/6,7/6$), tacnode ($3/4,5/4$) and $D_5$ ($5/8,\ldots$) do not lie in $\frac13\Z$, and on $\mathbb{CP}^1\times\mathbb{CP}^1$ ($s\in\frac12\Z$) the spectral numbers of $A_k$ have no half-integer other than $1$, so no anomaly appears.

\tV\ Unequal residues: for a triple point with module $\prod h_i^{s_i}$ ($h=zw(z-w)$), in 533 random $(s_1,s_2,s_3)\in(\frac16\Z\cap[-2,2])^3$ ($s_i\ne1$) the stalk is $\C^4[-2]\oplus R$, where $R=H^*(U_{\rm loc};L_s)$ if $s_1+s_2+s_3<2$ and $0$ otherwise. Whereas for a double point the threshold is per branch, the threshold at a triple point is on the \emph{sum} of the residues. \tK\ General principle: the threshold is the ratio of the weight of the module to the logarithmic canonical weight with respect to the symmetric torus action on the germ.

\begin{conj}[Refined Riemann--Roch (equal residues, isolated singularities)]\tK
Let $X$ be a compact surface, $\sigma\in H^0(-K_X)$ with zero curve $D$ reduced with only isolated singularities, $L$ a Poisson module with $c_1(L)=sK_X$, $L_s$ the rank-one local system on $U$ (monodromy $e^{2\pi is}$ around each branch), and $\mathcal F_s$ the constructible complex which is $L_s$ on $U$ and at $p\in D$ is the stalk of $Rj_*L_s$ if $s<\mathrm{lct}_p$ and $0$ if $s>\mathrm{lct}_p$. (a) For $s\notin\mathrm{Sp}(D)$ (the spectral numbers of all singularities of $D$, together with $1$): $\HH^*=H^*(X;\mathcal F_s)\oplus\bigoplus_p\C^{\mu_p}[-2]$. (b) The Euler characteristic is always $\chi(X)$ ($\chi(D)=\sum\mu_p$, from $\delta$-invariants \tP). (c) For $s\ll0$, $\mathcal F_s=Rj_*L_s$, and $\HH^*=H^*(U;L_s)\oplus\bigoplus\C^{\mu_p}[-2]$ depends only on $s\bmod\Z$ and vanishes in degrees $\ge3$. (d) For $s\in\mathrm{Sp}(D)$, rank-one cancellations can occur.
\tV\ On $\mathbb{CP}^2$ (smooth, nodal, cuspidal, conic plus secant, conic plus tangent, triangle, triple point) and $\mathbb{CP}^1\times\mathbb{CP}^1$ (4 examples), for $m=-15..15$ ($-9..8$), no contradiction except in (d). Not proved: the stalk results are only for quasi-homogeneous germs, and the global identification $\mathbb H(C_m)\simeq(\mathcal F_s+\text{skyscraper})$ is a fit. For non-reduced $D$ (double line, double conic) the stalks are infinite-dimensional and the global dimensions grow linearly in $m$ (slope $\deg\mathcal O(1)|_C$: $1$ for a double line on $\mathbb{CP}^2$, $2$ for a double conic on $\mathbb{CP}^1\times\mathbb{CP}^1$).
\end{conj}


%--------------------------------------------------------------
\section{Comparison with earlier work}\label{sec:compare}
%--------------------------------------------------------------
\subsection{Hodge theory (formal coefficients)}
Gualtieri proved a generalized Hodge star, the Born--Infeld inner product, Laplacian equalities, the Hodge decomposition and the $dd^{\Jc}$-lemma on generalized K\"ahler manifolds (\tC\ Gualtieri, \emph{Generalized geometry and the Hodge decomposition}, arXiv:math/0409093; Gualtieri, \emph{Generalized K\"ahler geometry}, Comm.\ Math.\ Phys.\ 331 (2014) 297--331, arXiv:1007.3485; Cavalcanti, \emph{New aspects of the $dd^c$-lemma}, arXiv:math/0501406, and J.\ Geom.\ Phys.\ 57 (2006) 121--132. Theorem numbers are omitted because these were not read closely \tQ). These are theories with trivial or $H$-twisted coefficients; a Bochner--Kodaira--Nakano identity with line-bundle (curved) coefficients has not been written, as far as surveyed. The results of this note can be positioned as filling this gap for $\beta$-transforms, not by Hodge theory but by comparing logarithmic complexes with constructible sheaves.

\subsection{Pal--Poddar}
\tC\ Pal--Poddar, \emph{Strong generalized holomorphic principal bundles}, arXiv:2404.18113, treat compact generalized complex manifolds of regular type (constant type $k$) whose symplectic leaf space is a K\"ahler orbifold, and give Hodge theory and Serre duality for strong generalized holomorphic (SGH) bundles (transversal Dolbeault isomorphism Thm.~8.9, Cor.~8.10 (numbers \tQ); Hodge decomposition and Serre duality Thm.~11.2, 11.4, 11.6, 11.7), K\"ahler identities (Prop.~11.8, Lemma~11.9; numbers \tQ), and Kodaira-type and Serre-type vanishing theorems for positive SGH line bundles (Def.~11.10, 11.11; Thm.~11.12). The hypotheses of regular type and orbifold leaf space are used to guarantee that the foliation is a smooth leaf bundle and to guarantee transversal Hodge theory (ellipticity). Differences from the present note:
\begin{itemize}
\item The type-jumping case (the type jumps on the degeneracy divisor of $\sigma$) is outside the framework of Pal--Poddar; examples such as $\mathbb{CP}^2$ in this note are of this kind.
\item SGH bundles form a proper subclass of generalized holomorphic bundles (flat along the leaves and holomorphic transversally). The Poisson modules of this note are general GH bundles on a $\beta$-transformed structure.
\item The statement and proof structure of Thm.~11.12 were checked by re-reading \S11: $H^q(M,(\mathcal G^*M)^p\otimes E)=0$ for $p+q>k$, and $H^q(M,E'\otimes E^m)=0$ for $m\ge m_0$, $q>0$. The hypotheses (compactness, a K\"ahler orbifold leaf space, positivity of $E$) are used for transversal Hodge theory and for positivity of the curvature term of a Bochner--Kodaira--Nakano type identity.
\end{itemize}
Doubts about the original text noticed in this survey (likely misprints or omissions in the source, but unverified \tQ): (i) the statement carries a factor $(i/2\pi)$ while the proof has no $1/2\pi$ (harmless); (ii) the sign discrepancy about $gc_1$ comes from the differing definitions of the characteristic polynomial in (9.9) and of $\det(I\pm tA/2\pi i)$ in \S10.1; (iii) ``a positive class gives a metric of positive curvature'' does not follow from Def.~11.11 and is assumed in the proof of Thm.~11.12; (iv) Def.~11.11 does not state compactness or the K\"ahler condition; (v) the proof of Lemma~11.9(1) is a computation in pointwise local frames; (vi) $T^{1,0}$ versus $T^{1,0*}$ in Thm.~13.4(2)(3), and the binomial coefficient for $q=n$ in Example 13.10.

\subsection{The work of Tengzhou Hu}
\tC\ T.\ Hu, \emph{Kodaira vanishing theorems for K\"ahler Lie algebroids}, arXiv:2403.10876;\\ T.\ Hu, \emph{Lie groupoid Riemann--Roch--Hirzebruch theorem and applications}, arXiv:2403.13424.
\begin{itemize}
\item \emph{Setting} (2403.13424): a Lie groupoid $G\rightrightarrows M$ (Def.~1; hypotheses of properness or $s$-connectedness are not stated) and its Lie algebroid $\mathfrak g$. A ``complex Lie algebroid'' (Def.~2) is a \emph{real} Lie algebroid $\mathfrak g$ with an integrable complex structure $J$ in the fibre direction (this is different from the complex Lie algebroid $L$ as a subbundle of $\bT_\C$). Kähler Lie algebroids are Def.~11, positive bundles Def.~12.
\item \emph{Claims} (2403.13424): Thm.~4 (Riemann--Roch) states, for $M$ compact, $\dbar_{\mathfrak g}$ elliptic (claimed, but no proof was found) and $\Omega$ a $\mathfrak g$-invariant density (unimodularity),
\[\mathrm{Ind}_\Omega(\dbar_{\mathfrak g}+\dbar_{\mathfrak g}^*)=(2\pi i)^{-k}\int_M(-1)^{\mathrm{rk}\,\mathfrak g}\bigl\langle\ch^{\mathfrak g}(\mathfrak g^{*p,0})\ch^{\mathfrak g}(E)\Td^{\mathfrak g}_\C(\mathfrak g),\Omega_{\mathfrak g}\bigr\rangle.\]
Thm.~7 (vanishing): if $E$ is positive and $p+q>n$, the \emph{kernel} of the $\mathfrak g$-Dolbeault Laplacian vanishes (vanishing of cohomology does not follow: for a general Lie algebroid $\mathfrak g$ the Laplacian is not necessarily elliptic and Hodge theory is unavailable). In 2403.10876, a K\"ahler Lie algebroid is Def.~2, the K\"ahler identities are Thm.~1 and Prop.~4, the BKN identity is Lemma~4, and the main vanishing theorem is in the introduction.
\item \emph{The generalized Dolbeault complex is outside the scope}\ \tP\ (derived; comparison with the source only to the extent indicated below): $L$ is a complex Lie algebroid without $J$, and $\wedge L^*$ in general has no $(p,q)$ bigrading. For a holomorphic Poisson structure $\sigma$ the real Lie algebroid $A_{\mathbb R}$ falls within Hu's setting, but its Dolbeault complex is $\wedge\mathfrak g^{*0,q}$ built from $\bar\sigma$, which differs from the complex $(T^{0,1}\bowtie A^{1,0})$ of this note.
\item \emph{Does the index see $m$?}\ \tK: Hu's index is a groupoid index, and is expected to depend polynomially on $m$, roughly like $\int\ch(E)\Td(T^*)\ch(\wedge^pT)$, whereas the total Euler characteristic of this note is independent of $m$ (Prop.~\ref{prop:euler}). Moreover the residue $-m/3$, which governs the periodic jumps, is not part of Hu's data. Hence Hu's theorem is not ``Riemann--Roch for generalized Dolbeault cohomology''.
\item \emph{Inconsistencies in the source}: 2403.13424 (obtained through summaries): $n$ is undefined in Thm.~7; the sign constants differ among Thm.~5, Prop.~1 and Thm.~6; ``unimodular'' is defined in two ways. 2403.10876 (checked by reading further and part of the appendix \tC): Lemma~5 gives $[L,\Lambda]=(p+q-n)\,\mathrm{id}$, but the proof of Thm.~2 writes $\Delta_{\bar\partial_{\mathfrak g}}=\Delta^{1,0}-[L,\Lambda]=\Delta^{1,0}+(p+q-n)\,\mathrm{id}$, and the sign in the last equality does not match ($-(p+q-n)=n-p-q$). Moreover the substitution $\nabla^2=\sqrt{-1}L$ means $\sqrt{-1}R=-L$, opposite to the definition of ``positive'' ($\sqrt{-1}R>0$, Def.~13). The statement itself (the kernel vanishes for positive $E$ when $p+q>n$) agrees with classical Kodaira--Nakano; so the printed proof has a sign slip, not the claim. Misprints in Cor.~1 (arrows and the symbols $s/t$). The appendix (Lemmas 6--10, the computation of $[\bar\partial_{\mathfrak g},\iota_\omega]$) could not be verified because the indices are hard to read. The reference list and publication data are unchecked. \tQ
\item \emph{Correction to the earlier version of this note (preliminary survey)}: the earlier version gave the author as ``S.~Hu'' and a different title; both were wrong, and the above is correct.
\end{itemize}

\subsection{Integrability and groupoids}
\tC\ Gualtieri--Li (arXiv:1206.3674, IMRN 2014) classify Hausdorff symplectic groupoids of log-symplectic manifolds and show that the Lie algebroid $T_X(-\log D)$ of logarithmic vector fields tangent to a smooth $D$ is integrable. Bailey--Gualtieri (arXiv:1611.03850) show the equivalence between integration of a generalized complex structure to a weak holomorphic symplectic groupoid and integrability of the real Poisson structure (Thm.~7.1: integrability of the underlying real Poisson structure is necessary and sufficient; Thm.~7.1 itself says nothing about Hausdorffness or $s$-connectedness, and the construction Thm.~7.6 uses an $s$-connected, $s$-simply connected integration). Laurent-Gengoux--Sti\'enon--Xu (arXiv:0803.2031): Thm.~3.17, 3.22. Crainic--Fernandes (arXiv:math/0105033): Thm.~4.1, Cor.~5.2. These provide groupoids as input for index theory, but do not treat integration of the complex Lie algebroid $L$ itself.

%--------------------------------------------------------------
\section{Open problems and next steps}\label{sec:open}
%--------------------------------------------------------------
\begin{enumerate}
\item \textbf{The treasure not yet touched}: a generalization of the Kodaira embedding theorem. Since $\HH^0$ consists of Casimir sections (Observation~\ref{obs:casimir}), it must be formulated not as an embedding of $M$ itself but as a map to the quotient separated by Casimir sections (e.g.\ the leaf space). The next concrete step is to study base points and separation properties of the linear system of Casimir sections in the $\mathbb{CP}^2$ example ($\HH^0=\C F^{m/3}$, i.e.\ 1-dimensional for $m\in3\Z_{\ge0}$).
\item \textbf{Generalization of Serre-type vanishing} (Conjecture~\ref{conj:A3}): a definition of a ``positive'' generalized holomorphic line bundle, and a construction of a filtration playing the role of the $E_1$ term for $\Jc$ that are not $\beta$-transforms. The comparison with Pal--Poddar Thm.~11.12 was done by re-reading \S11 (the groups differ; the transversal-degree reading is consistent).
\item \textbf{Formulation of Riemann--Roch}: instead of the total Euler characteristic, one must fix a refinement (the $E_1$ term of a filtration by the type of $\Jc$, a periodic quasi-polynomial, the indicial roots $H^*(\partial U)$). Proving Conjecture~\ref{conj:I3} requires an index theory for the complex Lie algebroid $L$ (unimodularity, and a derivation of Fredholmness $\Leftrightarrow r\notin\Z$). An APS-type extension to generalized complex manifolds with boundary is a still more distant problem.
\item \textbf{Beyond log-symplectic}: proofs for the local theory of non-nodal singularities (only quasi-homogeneous germs were computed), a global formula for unequal residues, non-reduced $D$ (infinite-dimensional stalks and linear growth), and non-toric examples in dimension $\ge4$ (see Observation~\ref{obs:arith}).
\item \textbf{Generalized K\"ahler side}: write down the Bochner--Kodaira--Nakano identity with line-bundle coefficients in the generalized K\"ahler framework. The interaction of the $H$-twist with the curvature term may be an obstacle \tK.
\item \textbf{Resolution of the unverified items of this note}: the \tQ\ items of \S\ref{sec:ledger}. In particular Mokri's theorem, a re-derivation of LGSX Lemma~4.22, a reading of Gualtieri 2011 \S4, and the remaining bibliographic data.
\end{enumerate}

\paragraph{Priorities for the next steps (author's proposal).}
(a) As a comparison case in which the type does not jump, compute the same ``Poincar\'e polynomial'' for the regular-type examples of Pal--Poddar and compare with the periodic jumps of this note. (b) For the $\mathbb{CP}^2$ example, use $\HH^0=\C F^{m/3}$ to study the geometry of the linear system of Casimir sections and look for the ``right form'' of a Kodaira-type theorem. (c) As a basis for an index-theoretic proof of Conjecture~\ref{conj:I3}, formulate rigorously the relation between the weighted $b$-complex on the real Lie algebroid $T_X(-\log D)$ and the residue $r$.


%--------------------------------------------------------------
\section{Status of the claims and extent of verification}\label{sec:ledger}
%--------------------------------------------------------------
\subsection{Status of each claim}
\begin{center}\footnotesize
\begin{tabular}{>{\raggedright\arraybackslash}p{4.0cm}>{\raggedright\arraybackslash}p{3.0cm}>{\raggedright\arraybackslash}p{8.4cm}}\toprule
Claim & Status & Dependencies / gaps\\\midrule
Theorem~\ref{thm:main1} (Poisson module $=$ log connection) & \tP & Any even dimension, simple normal crossings. The residue formula $c_1=-\sum r_i[D_i]$ is quoted from Esnault--Viehweg.\\
Theorem~\ref{thm:cp2} (table for $\mathbb{CP}^2$) & \tP\ \tV & Hand derivation and exact computation for $m=-15..15$ agree. Depends on the identification of complexes (Prop.~\ref{prop:ident}(b)).\\
Theorem~\ref{thm:main2} (main theorem for surfaces) & \tP & Prop.~\ref{prop:ident}(b) (LGSX Thm.~4.19 and Mokri's theorem), Deligne's comparison, the Laurent-tail computation of local cohomology (standard but not written out). Requires $r_i\ne1$.\\
Theorem~\ref{thm:B} (closed formula) & \tP\ \tV & Theorem~\ref{thm:main2} plus (H1)--(H4). (H4) ($H^1(U;L_\chi)=0$) was checked in 8 examples; in general it is a condition on $V^1(U)$. Agrees for 196 values.\\
Theorem~\ref{thm:A1} (Serre-type) & \tP & Only classical Serre (Hartshorne III.5.2) and the acyclic resolution lemma. Any dimension.\\
Theorem~\ref{thm:A2} (periodic stabilization) & \tP & Depends on Theorem~\ref{thm:main2} and its hypotheses ($X$ simply connected, $U$ affine).\\
Prop.~\ref{prop:euler} (Euler characteristic) & \tP & HRR only.\\
Theorem~\ref{thm:dual} (duality) & \tP\ \tV & Lemmas D and E machine-checked for $n\le4$. The agreement of the Koszul signs of the $\dbar$ part and the $d$ part is by hand only. Hodge theory for elliptic complexes is standard.\\
Theorem~\ref{thm:weight} (weight formula) & \tP\ \tV & Assumes $E_1^{a>0}=0$. Agreement with linear algebra in 60 examples, and in 3626 examples with generic $A$.\\
Obs.~\ref{obs:arith} (arithmetic dependence of $\sigma$) & \tV & A few examples in dimension 4. The explanation of the mechanism is a sketch.\\
Local theory in \S\ref{sec:local} and the triple-point explanation & \tV\ \tK & Quasi-homogeneous germs only. The global identification is a fit.\\
Conjectures~\ref{conj:A3}, \ref{conj:I3} & \tK & Neither definitions nor proofs are complete.\\
Generalization of the Kodaira embedding theorem & --- & Not treated (\S\ref{sec:open}).\\\bottomrule
\end{tabular}
\end{center}

\subsection{Unverified or delicate items (\tQ)}
\begin{itemize}
\item \textbf{Bibliographic data (from memory)}: the journal and pages of Goto, the IMRN volume/pages of Gualtieri--Li, the OUP data of Gualtieri ``Branes'', the volume/pages of Mokri, of Arapura, the journal of Lang--Jia--Liu, the IMRN Art.\ ID of LGSX, the volume/issue of Pal--Poddar. The numbers in Deligne LNM 163, Crainic--Fernandes, Fulton, Esnault--Viehweg. The numbers and contents of Thm.~8.9, Prop.~9.3, Ex.~10.18, Prop.~11.8, Lemma~11.9, Ex.~13.10 of Pal--Poddar. (Checked in v2: Hitchin Moscow Math.\ J.\ 12 (2012) 567--591, Pym Lett.\ Math.\ Phys.\ 108 (2018) 573--632, Hausel--Hunsicker--Mazzeo Duke 122 (2004) 485--548, Timmerscheidt J.\ reine angew.\ Math.\ 379 (1987) 152--171, title and authors of Pal--Poddar.)
\item \textbf{Gualtieri 2011 (arXiv:math/0703298)}: of roughly 160,000 characters, the first roughly 100,000 (up to the middle of \S3) were read closely; \S4 (the Darboux theorem etc.) was not read in full. In v2 the following were checked: type Def.~3.5, canonical bundle Def.~3.7, ellipticity Prop.~3.13/Cor.~3.14, Picard group Def.~3.18, Darboux Thm.~4.3 (regular points), regular points Def.~4.1; the type-jumping example is in \S4.1, $\rho=z_1+dz_1\wedge dz_2$ on $\mathbb R^4$.
\item \textbf{Hitchin, \emph{Poisson modules and generalized geometry} (arXiv:0905.3227, title and author checked)}: it only states in the introduction that GH bundles $=$ Poisson modules; it contains no $\Jc_\sigma$ notation, no identification of hypercohomology, no obstruction by the Atiyah class, no $\mathbb{CP}^2$ example and no Lie algebroid $L$ (according to the text obtained). This note instead cites Branes on Poisson varieties (arXiv:0710.2719) and LGSX. (v1 attributed this paper to Gualtieri and said ``confirmed''; that was wrong in both author and content.)
\item \textbf{$\mathrm{Pic}(\Jc)=\mathbb H^1(\mathcal L_{\log})$}: could not be checked in the text of any reference. This note does not use it. What was used is only the direct derivation of existence and uniqueness of Poisson module structures.
\item \textbf{LGSX Lemma~4.22, Mokri's theorem}: not re-derived. The standard facts of Hodge theory and elliptic complexes, and Esnault--Viehweg, Deligne, Fulton--Deligne, Timmerscheidt, Arapura and Crainic--Fernandes, are quoted, and their ranges of applicability were not checked against the sources.
\item \textbf{The two papers of Hu}: for 2403.10876 the first roughly 200,000 characters (about 500,000 in the full expanded HTML); for 2403.13424 only the overview and the statements of theorems.
\item \textbf{Pal--Poddar}: \S\S5--11, 13 were re-read (22 November 2026); the extraction went through summaries, so details of formulas are unverified.
\item \textbf{Limits of the computations}: exact linear algebra was done for a finite range of $m$, a finite number of curves (Fermat and Hesse cubics, etc.), and modulo primes ($32003$, $2^{61}-1$). Dependence on the choice of curve (random smooth curves), and whether reduction modulo $p$ represents the characteristic-0 result, were not verified.
\end{itemize}

\subsection{Record of the work}
All scripts use exact arithmetic in Python (standard library) and were saved in \texttt{cp2/} (\texttt{poisson.py, table.py, neg\_check.py, predict.py, sing\_table.py, p1p1.py, local\_stalk.py, local\_unequal.py, ncd\_cochain\_check.py, ncd\_local\_weights.py, ncd\_global\_check.py, serre\_rr\_check.py, serre\_dual\_check.py, duality\_ndim.py, duality\_ndim\_exact.py, toric\_hh.py, toric\_crit.py, toric\_generic\_check2.py, index\_route\_check.py}).

%--------------------------------------------------------------
\section{Revision notes (v1 $\to$ v2, 22 November 2026)}\label{sec:rev}
%--------------------------------------------------------------
Errors and updates found by re-checking the citations of v1 (arXiv pages, publisher pages, re-reading of the sources). No mathematical statement (theorem or proposition) has been changed.
\begin{enumerate}
\item \textbf{Author of arXiv:0905.3227}: \textbf{Hitchin}, not Gualtieri (\emph{Poisson modules and generalized geometry}). The bibliography and \S\ref{sec:ledger} are corrected.
\item \textbf{Numbering in Gualtieri 2011}: ellipticity is Prop.~3.13 and Cor.~3.14 (``Prop.~3.16'' of v1 does not exist); Def.~3.18 is the Picard group, and generalized holomorphic bundles are introduced in prose after Thm.~3.15; the generalized Darboux theorem is Thm.~4.3 (``4.35'' of v1 was wrong); the ``Examples 4.1--4.2'' of v1 do not exist, Def.~4.1 being the definition of a regular point.
\item \textbf{Pal--Poddar (arXiv:2404.18113)}: title \emph{Strong generalized holomorphic principal bundles}, authors D.~Pal and M.~Poddar; the ``author discrepancy'' of v1 did not reproduce. Theorem numbers organized (Hodge decomposition and Serre duality are Thm.~11.2, 11.4, 11.6, 11.7; vanishing is Thm.~11.12). The consistency of Thm.~11.12 with the transversal-degree reading of this note was checked by re-reading \S\S5--11, 13 (the groups are not identical). The cause of the sign doubts in the source (differing definitions of the characteristic polynomial) was identified.
\item \textbf{Lang--Jia--Liu}: arXiv:2102.07177, \emph{The Atiyah class of generalized holomorphic vector bundles}; GH bundles defined by local trivializations, a notion different from the SGH bundles of Pal--Poddar.
\item \textbf{Hu (Tengzhou Hu), 2403.10876}: a sign slip in the printed proof of the vanishing theorem was confirmed by reading further in the text (the statement agrees with the classical one).
\item \textbf{Bibliographic data fixed}: Hitchin (Moscow Math.\ J.\ 12 (2012) 567--591), Pym (Lett.\ Math.\ Phys.\ 108 (2018) 573--632), Hausel--Hunsicker--Mazzeo (Duke 122 (2004) 485--548), Timmerscheidt (J.\ reine angew.\ Math.\ 379 (1987) 152--171; the recollection in v1 was wrong), year and DOI of Mokri, arXiv numbers of Goto and Arapura.
\item \textbf{Bailey--Gualtieri Thm.~7.1}: confirmed that integrability of the real Poisson structure is necessary and sufficient; a remark on Hausdorffness and $s$-connectedness added.
\end{enumerate}
The \tQ\ items not declared checked in this section remain in \S\ref{sec:ledger}.


\begin{thebibliography}{99}\small
\bibitem{AB} M.\ Abouzaid, M.\ Boyarchenko, Local structure of generalized complex manifolds, J.\ Symplectic Geom.\ 4 (2006) 43--62, arXiv:math/0511049.
\bibitem{AN} Y.\ Akizuki, S.\ Nakano, Note on Kodaira--Spencer's proof of Lefschetz's theorems, Proc.\ Japan Acad.\ 30 (1954) 266--272.
\bibitem{ALN} B.\ Ammann, R.\ Lauter, V.\ Nistor, Pseudodifferential operators on manifolds with a Lie structure at infinity, arXiv:math/0304044.
\bibitem{Ara} D.\ Arapura, Geometry of cohomology support loci for local systems I, J.\ Algebraic Geom.\ 6 (1997) 563--597, arXiv:alg-geom/9612006 [volume/pages from memory \tQ].
\bibitem{AS} M.\ Atiyah, I.\ Singer, The index of elliptic operators III, Ann.\ Math.\ 87 (1968) 546--604.
\bibitem{BG} A.\ Bailey, M.\ Gualtieri, Integration of generalized complex structures, arXiv:1611.03850.
\bibitem{BLS} M.\ Braverman, Y.\ Loizides, Y.\ Song, Geometric quantization of $b$-symplectic manifolds, arXiv:1910.10016.
\bibitem{Cav} G.\ Cavalcanti, New aspects of the $dd^c$-lemma, arXiv:math/0501406; The decomposition of forms and cohomology of generalized complex manifolds, J.\ Geom.\ Phys.\ 57 (2006) 121--132.
\bibitem{CF} M.\ Crainic, R.\ L.\ Fernandes, Integrability of Lie brackets, Ann.\ Math.\ 157 (2003) 575--620, arXiv:math/0105033.
\bibitem{Del} P.\ Deligne, \'Equations diff\'erentielles \`a points singuliers r\'eguliers, Lecture Notes in Math.\ 163, Springer 1970.
\bibitem{DelFG} P.\ Deligne, Le groupe fondamental du compl\'ement d'une courbe plane n'ayant que des points doubles ordinaires est ab\'elien, S\'em.\ Bourbaki 543 (1979/80).
\bibitem{EV} H.\ Esnault, E.\ Viehweg, Logarithmic de Rham complexes and vanishing theorems, Invent.\ Math.\ 86 (1986) 161--194.
\bibitem{Ful} W.\ Fulton, On the fundamental group of the complement of a node curve, Ann.\ Math.\ 111 (1980) 407--409.
\bibitem{Goto} R.\ Goto, Unobstructed deformations of generalized complex structures induced by $C^\infty$ logarithmic symplectic structures, arXiv:1501.03398 [journal and pages unverified \tQ].
\bibitem{GH} P.\ Griffiths, J.\ Harris, Principles of Algebraic Geometry, Wiley 1978.
\bibitem{GuaTh} M.\ Gualtieri, Generalized complex geometry, Oxford DPhil thesis, arXiv:math/0401221.
\bibitem{GuaHodge} M.\ Gualtieri, Generalized geometry and the Hodge decomposition, arXiv:math/0409093.
\bibitem{Gua11} M.\ Gualtieri, Generalized complex geometry, Ann.\ Math.\ 174 (2011) 75--123, arXiv:math/0703298.
\bibitem{GuaK} M.\ Gualtieri, Generalized K\"ahler geometry, Comm.\ Math.\ Phys.\ 331 (2014) 297--331, arXiv:1007.3485.
\bibitem{GuaBr} M.\ Gualtieri, Branes on Poisson varieties, in: The Many Facets of Geometry, arXiv:0710.2719 (v2, July 2010; the arXiv page gives no journal data; The Many Facets of Geometry, Oxford Univ.\ Press 2010 is from memory \tQ).
\bibitem{GuaPM} N.\ Hitchin, Poisson modules and generalized geometry, arXiv:0905.3227 (2009) (title and author checked on the arXiv page; publication data unverified \tQ).
\bibitem{GL} M.\ Gualtieri, S.\ Li, Symplectic groupoids of log symplectic manifolds, IMRN 2014, arXiv:1206.3674.
\bibitem{Hart} R.\ Hartshorne, Algebraic Geometry, GTM 52, Springer 1977.
\bibitem{HHM} T.\ Hausel, E.\ Hunsicker, R.\ Mazzeo, Hodge cohomology of gravitational instantons, Duke Math.\ J.\ 122 (2004) 485--548, arXiv:math/0207169.
\bibitem{Hirz} F.\ Hirzebruch, Neue topologische Methoden in der algebraischen Geometrie, Springer 1956.
\bibitem{Hit03} N.\ Hitchin, Generalized Calabi--Yau manifolds, Q.\ J.\ Math.\ 54 (2003) 281--308, arXiv:math/0209099.
\bibitem{Hit12} N.\ Hitchin, Deformations of holomorphic Poisson manifolds, Moscow Math.\ J.\ 12 (2012), no.\ 3, 567--591, arXiv:1105.4775.
\bibitem{Hu1} T.\ Hu, Kodaira vanishing theorems for K\"ahler Lie algebroids, arXiv:2403.10876.
\bibitem{Hu2} T.\ Hu, Lie groupoid Riemann--Roch--Hirzebruch theorem and applications, arXiv:2403.13424.
\bibitem{Huy} D.\ Huybrechts, Complex Geometry: An Introduction, Springer 2005.
\bibitem{Kod53} K.\ Kodaira, On a differential-geometric method in the theory of analytic stacks, Proc.\ Nat.\ Acad.\ Sci.\ USA 39 (1953) 1268--1273.
\bibitem{Kod54} K.\ Kodaira, On K\"ahler varieties of restricted type, Ann.\ Math.\ 60 (1954) 28--48.
\bibitem{Lan} M.\ Lanius, Poisson cohomology of a class of log symplectic manifolds, arXiv:1605.03854.
\bibitem{LGSX1} C.\ Laurent-Gengoux, M.\ Sti\'enon, P.\ Xu, Holomorphic Poisson manifolds and holomorphic Lie algebroids, IMRN 2008, arXiv:0707.4253.
\bibitem{LGSX2} C.\ Laurent-Gengoux, M.\ Sti\'enon, P.\ Xu, Integration of holomorphic Lie algebroids, arXiv:0803.2031.
\bibitem{LJL} H.\ Lang, X.\ Jia, Z.\ J.\ Liu, The Atiyah class of generalized holomorphic vector bundles, arXiv:2102.07177; Differential Geom.\ Appl.\ 90 (2023) 102031 (journal data taken from the reference list of Pal--Poddar \tQ). GH bundles are defined by local trivializations; this is a different notion from the SGH bundles of Pal--Poddar.
\bibitem{MM} R.\ Mazzeo, R.\ Melrose, Pseudodifferential operators on manifolds with fibred boundaries, arXiv:math/9812120.
\bibitem{Mel} R.\ Melrose, The Atiyah--Patodi--Singer Index Theorem, A K Peters 1993.
\bibitem{Mok} T.\ Mokri, Matched pairs of Lie algebroids, Glasgow Math.\ J.\ 39 (1997) 167--181, doi:10.1017/S0017089500032055 [year and DOI checked on the publisher page, volume/pages from memory \tQ].
\bibitem{PP} D.\ Pal, M.\ Poddar, Strong generalized holomorphic principal bundles, arXiv:2404.18113; Ann.\ Sc.\ Norm.\ Super.\ Pisa Cl.\ Sci., doi:10.2422/2036-2145.202501\_026 (volume/issue \tQ).
\bibitem{PPT} M.\ Pflaum, H.\ Posthuma, X.\ Tang, The index of geometric operators on Lie groupoids, Indag.\ Math.\ 25 (2014) 1--18, arXiv:1308.0236.
\bibitem{Pym} B.\ Pym, Constructions and classifications of projective Poisson varieties, Lett.\ Math.\ Phys.\ 108 (2018), no.\ 3, 573--632, doi:10.1007/s11005-017-0984-5.
\bibitem{QS} Y.\ Qiao, Y.\ So, Deformation and K-theoretic index formulae on boundary groupoids, arXiv:2304.13345.
\bibitem{Serre} J.-P.\ Serre, Faisceaux alg\'ebriques coh\'erents, Ann.\ Math.\ 61 (1955) 197--278.
\bibitem{Tim} K.\ Timmerscheidt, Mixed Hodge theory for unitary local systems, J.\ reine angew.\ Math.\ 379 (1987) 152--171.
\end{thebibliography}


\end{document}
